\chapter{Линейная алгебра}
\label{chap:linear-algebra}

Линейная алгебра описывает объекты, которые можно складывать и умножать на числа, а также отображения, сохраняющие эти операции. Этот язык нужен не только для решения систем уравнений. Векторы представляют параметры, признаки, состояния и направления изменения; матрицы задают преобразования; базис определяет способ записи объекта числами. Поэтому выбор координат, размерность и линейная зависимость непосредственно влияют на устройство вычислительной модели.

В этой главе конечномерная теория излагается с опорой на курс Э.~Б.~Винберга. Доказательства приводятся лишь там, где они объясняют устройство понятия или метода. Основное внимание уделяется тому, как распознать объект, как выполнить вычисление и как интерпретировать результат.

\section{Векторные пространства, базис и координаты}
\label{sec:vector-spaces-bases}

\subsection{Поле скаляров и векторное пространство}

Пусть \(K\) --- поле. Это означает, что элементы \(K\) можно складывать, вычитать, умножать и делить на любой ненулевой элемент, причём обычные правила арифметики сохраняются. В основных приложениях далее используются поля \(\R\) и \(\C\).

\begin{definition}
Векторным пространством над полем \(K\) называется множество \(V\), на котором определены сложение векторов
\[
V\times V\to V,
\qquad
(u,v)\mapsto u+v,
\]
и умножение вектора на скаляр
\[
K\times V\to V,
\qquad
(\lambda,v)\mapsto \lambda v,
\]
такие, что для любых \(u,v,w\in V\) и \(\lambda,\mu\in K\) выполняются соотношения
\[
\begin{aligned}
u+v&=v+u,
&
(u+v)+w&=u+(v+w),
\\
u+0&=u,
&
u+(-u)&=0,
\\
\lambda(u+v)&=\lambda u+\lambda v,
&
(\lambda+\mu)u&=\lambda u+\mu u,
\\
(\lambda\mu)u&=\lambda(\mu u),
&
1u&=u.
\end{aligned}
\]
Элементы \(V\) называются векторами, а элементы \(K\) --- скалярами.
\end{definition}

Определение не требует, чтобы вектор был геометрической стрелкой. Вектором может быть строка чисел, матрица, функция, многочлен или поле значений на сетке. Существенно только то, что сложение и умножение на скаляр удовлетворяют указанным правилам.

\begin{example}
Пространство \(K^n\) состоит из строк
\[
x=(x_1,\dots,x_n),
\qquad x_i\in K.
\]
Операции выполняются покоординатно:
\[
(x_1,\dots,x_n)+(y_1,\dots,y_n)
=(x_1+y_1,\dots,x_n+y_n),
\]
\[
\lambda(x_1,\dots,x_n)
=(\lambda x_1,\dots,\lambda x_n).
\]
Именно эта модель чаще всего реализуется в программе как одномерный массив чисел.
\end{example}

\begin{example}
Множество \(C[a,b]\) всех непрерывных вещественных функций на отрезке \([a,b]\) является векторным пространством над \(\R\). Сложение и умножение на число выполняются поточечно. Здесь один вектор содержит не конечный список координат, а целую функцию.
\end{example}

\subsection{Подпространства}

Внутри векторного пространства часто рассматривается множество допустимых состояний, удовлетворяющих ограничениям. Чтобы на этом множестве сохранилась линейная структура, оно должно быть замкнуто относительно линейных операций.

\begin{definition}
Подмножество \(U\subseteq V\) называется линейным подпространством пространства \(V\), если \(0\in U\) и для любых \(u,v\in U\), \(\lambda\in K\) выполнено
\[
u+v\in U,
\qquad
\lambda u\in U.
\]
\end{definition}

Эти условия можно объединить в один практический критерий.

\begin{criterion}[критерий подпространства]
Непустое подмножество \(U\subseteq V\) является подпространством тогда и только тогда, когда
\[
\alpha u+\beta v\in U
\]
для любых \(u,v\in U\) и любых \(\alpha,\beta\in K\).
\end{criterion}

Критерий проверяет, что любая линейная смесь допустимых объектов остаётся допустимой. Например, множество решений однородной системы \(Ax=0\) является подпространством. Множество решений неоднородной системы \(Ax=b\) при \(b\neq 0\) обычно подпространством не является, потому что не содержит нулевой вектор.

\subsection{Линейные комбинации и линейная оболочка}

\begin{definition}
Пусть \(v_1,\dots,v_m\in V\). Вектор вида
\[
\lambda_1v_1+\cdots+\lambda_mv_m,
\qquad
\lambda_1,\dots,\lambda_m\in K,
\]
называется линейной комбинацией векторов \(v_1,\dots,v_m\).
\end{definition}

Коэффициенты показывают, в каких долях и с какими знаками используются исходные направления. Все векторы, которые можно получить такими комбинациями, образуют подпространство.

\begin{definition}
Линейной оболочкой системы \(v_1,\dots,v_m\) называется множество
\[
\operatorname{span}(v_1,\dots,v_m)
=
\left\{
\lambda_1v_1+\cdots+\lambda_mv_m
\;\middle|\;
\lambda_1,\dots,\lambda_m\in K
\right\}.
\]
\end{definition}

Линейная оболочка --- это наименьшее подпространство, содержащее все \(v_i\). В вычислительной задаче она описывает все состояния, которые могут быть построены из заданного набора направлений.

\begin{example}
Пусть
\[
v_1=(1,0,1),
\qquad
v_2=(0,1,1).
\]
Тогда
\[
\alpha v_1+\beta v_2=(\alpha,\beta,\alpha+\beta),
\]
поэтому
\[
\operatorname{span}(v_1,v_2)
=
\left\{(x_1,x_2,x_3)\in\R^3
\;\middle|\;
x_3=x_1+x_2\right\}.
\]
Два вектора задают плоскость, проходящую через начало координат.
\end{example}

\subsection{Линейная зависимость}

\begin{definition}
Система \(v_1,\dots,v_m\) называется линейно зависимой, если существуют коэффициенты \(\lambda_1,\dots,\lambda_m\), не все равные нулю, такие, что
\[
\lambda_1v_1+\cdots+\lambda_mv_m=0.
\]
Если такое равенство возможно только при
\[
\lambda_1=\cdots=\lambda_m=0,
\]
то система называется линейно независимой.
\end{definition}

Линейная зависимость означает избыточность: хотя бы одно направление не добавляет новых возможностей, потому что выражается через остальные.

\begin{proposition}
Система \(v_1,\dots,v_m\) при \(m>1\) линейно зависима тогда и только тогда, когда хотя бы один её вектор выражается линейной комбинацией остальных.
\end{proposition}

\begin{proofidea}
Если, например,
\[
v_1=\mu_2v_2+\cdots+\mu_mv_m,
\]
то перенос всех слагаемых в одну часть даёт нетривиальную комбинацию, равную нулю. Обратно, из зависимости
\[
\lambda_1v_1+\cdots+\lambda_mv_m=0
\]
можно выбрать ненулевой коэффициент \(\lambda_j\) и выразить \(v_j\) через остальные векторы делением на \(\lambda_j\).
\end{proofidea}

В численной работе зависимость обнаруживается решением однородной системы. Если векторы записаны столбцами матрицы
\[
A=
\begin{pmatrix}
| & & | \\
v_1 & \cdots & v_m\\
| & & |
\end{pmatrix},
\]
то они линейно зависимы тогда и только тогда, когда система
\[
A\lambda=0
\]
имеет ненулевое решение.

\subsection{Базис, координаты и размерность}

\begin{definition}
Базисом векторного пространства \(V\) называется линейно независимая система
\[
e_1,\dots,e_n,
\]
линейная оболочка которой совпадает с \(V\).
\end{definition}

Два требования в определении выполняют разные роли. Порождаемость гарантирует, что базисных векторов достаточно для записи любого элемента пространства. Линейная независимость гарантирует единственность такой записи.

\begin{proposition}
Система \(e_1,\dots,e_n\) является базисом пространства \(V\) тогда и только тогда, когда каждый \(v\in V\) единственным образом представляется в виде
\[
v=x_1e_1+\cdots+x_ne_n.
\]
\end{proposition}

\begin{definition}
Числа \(x_1,\dots,x_n\) в единственном разложении
\[
v=x_1e_1+\cdots+x_ne_n
\]
называются координатами вектора \(v\) в базисе \(e_1,\dots,e_n\). Столбец координат обозначим через
\[
[v]_e=
\begin{pmatrix}
x_1\\
\vdots\\
x_n
\end{pmatrix}.
\]
\end{definition}

Вектор и столбец его координат --- не одно и то же. Вектор является абстрактным объектом пространства, а координатный столбец зависит от выбранного базиса. Смена базиса меняет числа, но не сам вектор.

\begin{theorem}
Все базисы конечномерного векторного пространства содержат одинаковое число векторов.
\end{theorem}

\begin{definition}
Размерностью конечномерного пространства \(V\) называется число векторов в любом его базисе. Используется обозначение
\[
\dim V=n.
\]
\end{definition}

В пространстве размерности \(n\) любая система из более чем \(n\) векторов линейно зависима. Любая линейно независимая система из ровно \(n\) векторов уже является базисом.

\begin{example}
Рассмотрим в \(\R^3\) векторы
\[
e_1=(1,1,0),
\qquad
e_2=(1,0,1),
\qquad
e_3=(0,1,1).
\]
Матрица с этими столбцами имеет вид
\[
A=
\begin{pmatrix}
1&1&0\\
1&0&1\\
0&1&1
\end{pmatrix}.
\]
Её определитель равен \(-2\), поэтому столбцы линейно независимы и образуют базис \(\R^3\). Для вектора \(v=(3,2,1)\) координаты находятся из системы
\[
A[v]_e=v.
\]
Решение даёт
\[
[v]_e=
\begin{pmatrix}
2\\1\\0
\end{pmatrix},
\qquad
v=2e_1+e_2.
\]
Вычислительно нахождение координат сводится к решению линейной системы.
\end{example}

\subsection{Смена базиса}

Пусть \(e=(e_1,\dots,e_n)\) и \(e'=(e'_1,\dots,e'_n)\) --- два базиса одного пространства. Выразим новый базис через старый:
\[
e'_j=\sum_{i=1}^n c_{ij}e_i.
\]
Матрица
\[
C=(c_{ij})
\]
называется матрицей перехода от базиса \(e\) к базису \(e'\). Её \(j\)-й столбец состоит из координат вектора \(e'_j\) в старом базисе.

Если один и тот же вектор имеет координатные столбцы \([v]_e\) и \([v]_{e'}\), то
\[
[v]_e=C[v]_{e'},
\qquad
[v]_{e'}=C^{-1}[v]_e.
\]
Первая формула собирает вектор из новых координат через столбцы нового базиса, записанные в старых координатах. Вторая формула переводит известные старые координаты в новые.

\begin{example}
В \(\R^2\) возьмём стандартный базис
\[
e_1=(1,0),
\qquad
e_2=(0,1)
\]
и новый базис
\[
e'_1=(1,1),
\qquad
e'_2=(1,-1).
\]
Матрица перехода равна
\[
C=
\begin{pmatrix}
1&1\\
1&-1
\end{pmatrix}.
\]
Для вектора \(v=(4,2)\) имеем
\[
[v]_{e'}=C^{-1}[v]_e
=
\frac12
\begin{pmatrix}
1&1\\
1&-1
\end{pmatrix}
\begin{pmatrix}
4\\2
\end{pmatrix}
=
\begin{pmatrix}
3\\1
\end{pmatrix}.
\]
Следовательно,
\[
v=3e'_1+e'_2.
\]
\end{example}

\paragraph{Связь с машинным обучением.}

Вектор признаков, скрытое представление нейросети и набор коэффициентов геометрической модели являются координатами относительно некоторой системы направлений. Замена базиса может разделить связанные факторы, уменьшить число существенных координат или сделать оператор почти диагональным. Позднее именно эта идея появится в спектральном разложении, методе главных компонент и сингулярном разложении.

\subsection*{Задачи}

\begin{problem}
\textit{Задача Винберга.}
Пусть \(V\) --- \(n\)-мерное пространство над полем \(\F_q\), содержащим \(q\) элементов. Найти число элементов пространства \(V\).
\end{problem}

\begin{proof}[Решение]
Выберем базис \(e_1,\dots,e_n\). Каждый вектор единственным образом задаётся координатами
\[
(x_1,\dots,x_n),
\qquad x_i\in\F_q.
\]
Для каждой из \(n\) координат имеется \(q\) независимых вариантов. Поэтому
\[
|V|=q^n.
\]
Смысл ответа состоит в том, что выбор базиса устанавливает взаимно однозначное соответствие между абстрактным пространством \(V\) и множеством координатных строк \(\F_q^n\).
\end{proof}

\begin{problem}
\textit{Задача Винберга.}
Показать, что пространство непрерывных вещественных функций на невырожденном промежутке бесконечномерно.
\end{problem}

\begin{proof}[Решение]
Рассмотрим функции
\[
1,x,x^2,\dots,x^m.
\]
Они непрерывны. Предположим, что существует линейная зависимость
\[
a_0+a_1x+\cdots+a_mx^m=0
\]
для всех \(x\) из промежутка. Левая часть является многочленом, имеющим бесконечно много корней. Такой многочлен может быть только нулевым, поэтому
\[
a_0=a_1=\cdots=a_m=0.
\]
Следовательно, для любого \(m\) в пространстве можно найти \(m+1\) линейно независимых функций. Конечного базиса существовать не может.
\end{proof}

\begin{problem}
\textit{Практическая задача.}
В пространстве \(\R^3\) задано множество
\[
U=\left\{x\in\R^3\;\middle|\;x_1+x_2+x_3=0\right\}.
\]
Проверить, что \(U\) является подпространством, построить его базис и найти координаты вектора
\[
v=(2,-1,-1)
\]
в найденном базисе.
\end{problem}

\begin{proof}[Решение]
Условие однородно. Если \(x,y\in U\) и \(\alpha,\beta\in\R\), то
\[
(\alpha x_1+\beta y_1)
+(\alpha x_2+\beta y_2)
+(\alpha x_3+\beta y_3)
=
\alpha\cdot 0+\beta\cdot 0=0.
\]
Значит, \(\alpha x+\beta y\in U\), и \(U\) является подпространством.

Из уравнения \(x_1+x_2+x_3=0\) получаем
\[
x_3=-x_1-x_2.
\]
Поэтому любой \(x\in U\) представляется как
\[
x=(x_1,x_2,-x_1-x_2)
=x_1(1,0,-1)+x_2(0,1,-1).
\]
Векторы
\[
b_1=(1,0,-1),
\qquad
b_2=(0,1,-1)
\]
линейно независимы и порождают \(U\), следовательно, образуют базис. Для заданного вектора
\[
v=2b_1-b_2,
\]
так что
\[
[v]_b=
\begin{pmatrix}
2\\-1
\end{pmatrix}.
\]
Размерность пространства допустимых состояний уменьшилась с трёх до двух: одно независимое линейное ограничение исключило одну степень свободы.
\end{proof}

\begin{problem}
\textit{Практическая задача.}
Даны векторы
\[
v_1=(1,0,1),
\qquad
v_2=(0,1,1),
\qquad
v_3=(1,1,2).
\]
Определить, образуют ли они базис \(\R^3\), и описать их линейную оболочку.
\end{problem}

\begin{proof}[Решение]
Непосредственно видно, что
\[
v_3=v_1+v_2.
\]
Следовательно, система линейно зависима и базисом \(\R^3\) не является. Третий вектор не добавляет нового направления, поэтому
\[
\operatorname{span}(v_1,v_2,v_3)
=
\operatorname{span}(v_1,v_2).
\]
Как было вычислено выше,
\[
\operatorname{span}(v_1,v_2)
=
\left\{(x_1,x_2,x_3)\in\R^3
\;\middle|\;
x_3=x_1+x_2\right\}.
\]
Это двумерная плоскость. В прикладной модели наличие \(v_3\) означает дублирование уже имеющейся степени свободы.
\end{proof}

\paragraph{Связь с дальнейшими разделами.}

Понятия линейной оболочки, базиса и размерности будут использоваться при изучении ранга матрицы, ядра и образа линейного отображения. Смена базиса станет основным инструментом для упрощения матриц операторов, а идея удаления избыточных направлений перейдёт в методы понижения размерности.

\section{Системы линейных уравнений и метод Гаусса}
\label{sec:linear-systems-gauss}

Система линейных уравнений возникает, когда неизвестный объект должен одновременно удовлетворять нескольким линейным ограничениям. В геометрии это может быть точка пересечения плоскостей, в механике --- состояние равновесия, в обработке данных --- набор параметров, точно воспроизводящий наблюдения. Метод Гаусса приводит такую систему к форме, из которой сразу видны совместность, число степеней свободы и все решения.

\subsection{Матрица и матричная запись системы}

\begin{definition}
Матрицей размера \(m\times n\) над полем \(K\) называется прямоугольная таблица
\[
A=(a_{ij})
=
\begin{pmatrix}
a_{11}&\cdots&a_{1n}\\
\vdots&&\vdots\\
a_{m1}&\cdots&a_{mn}
\end{pmatrix},
\qquad
a_{ij}\in K.
\]
Число \(a_{ij}\) находится в \(i\)-й строке и \(j\)-м столбце.
\end{definition}

Матрица действует на столбец \(x\in K^n\) по правилу
\[
(Ax)_i=\sum_{j=1}^n a_{ij}x_j.
\]
Поэтому система \(m\) линейных уравнений с \(n\) неизвестными
\[
\begin{aligned}
a_{11}x_1+\cdots+a_{1n}x_n&=b_1,\\
&\vdots\\
a_{m1}x_1+\cdots+a_{mn}x_n&=b_m
\end{aligned}
\]
записывается одним уравнением
\[
Ax=b,
\]
где
\[
x=
\begin{pmatrix}
x_1\\ \vdots\\ x_n
\end{pmatrix},
\qquad
b=
\begin{pmatrix}
b_1\\ \vdots\\ b_m
\end{pmatrix}.
\]

Эта запись имеет важную геометрическую интерпретацию. Если столбцы матрицы \(A\) обозначить через \(a_1,\dots,a_n\), то
\[
Ax=x_1a_1+\cdots+x_na_n.
\]
Следовательно, решить \(Ax=b\) означает выяснить, можно ли представить \(b\) линейной комбинацией столбцов \(A\), и найти коэффициенты этой комбинации.

\begin{definition}
Матрицей коэффициентов системы называется \(A\), а расширенной матрицей --- таблица
\[
[A\mid b]
=
\left(
\begin{array}{ccc|c}
a_{11}&\cdots&a_{1n}&b_1\\
\vdots&&\vdots&\vdots\\
a_{m1}&\cdots&a_{mn}&b_m
\end{array}
\right).
\]
\end{definition}

Вертикальная черта только отделяет коэффициенты от правой части. Все преобразования строк должны одновременно применяться и к \(A\), и к \(b\).

\subsection{Эквивалентные системы и элементарные преобразования}

\begin{definition}
Две системы называются эквивалентными, если множества их решений совпадают.
\end{definition}

Метод Гаусса использует три преобразования строк расширенной матрицы:
перестановку двух строк, умножение строки на ненулевой скаляр и прибавление к одной строке другой строки, умноженной на скаляр. Каждое такое действие соответствует обратимому преобразованию уравнений, поэтому не меняет множество решений.

\begin{proposition}
Элементарные преобразования строк переводят систему линейных уравнений в эквивалентную систему.
\end{proposition}

\begin{proofidea}
У каждого элементарного преобразования есть обратное. Перестановка отменяется той же перестановкой; умножение на \(\lambda\neq 0\) --- умножением на \(\lambda^{-1}\); прибавление \(\lambda\)-кратной строки --- прибавлением той же строки с коэффициентом \(-\lambda\). Поэтому из новой системы можно восстановить исходную без потери или появления решений.
\end{proofidea}

В матричном языке элементарное преобразование строк равносильно умножению слева на обратимую элементарную матрицу \(E\):
\[
Ax=b
\quad\Longleftrightarrow\quad
EAx=Eb.
\]
Это объясняет, почему преобразуется вся строка расширенной матрицы, включая свободный член.

\subsection{Ступенчатая форма}

\begin{definition}
Ведущим элементом ненулевой строки называется её первый ненулевой элемент. Матрица имеет ступенчатый вид, если ведущий элемент каждой следующей ненулевой строки расположен правее ведущего элемента предыдущей строки, а все нулевые строки находятся внизу.
\end{definition}

Позиции ведущих элементов называются ведущими позициями, или пивотами. Столбцы неизвестных, содержащие пивоты, соответствуют главным переменным. Остальные неизвестные называются свободными.

\begin{theorem}
Любую матрицу можно привести элементарными преобразованиями строк к ступенчатому виду.
\end{theorem}

Метод строится последовательно. В первом ещё не обработанном столбце выбирают ненулевой элемент, при необходимости переставляют строки и с его помощью зануляют элементы ниже. Затем переходят к подматрице, расположенной правее и ниже найденного пивота. Эта часть называется прямым ходом метода Гаусса.

После прямого хода решения находятся снизу вверх. Можно ограничиться обратной подстановкой либо дополнительно занулить элементы над каждым пивотом и нормировать пивоты до единицы. Полученная форма называется приведённой ступенчатой формой. Этот этап называется обратным ходом.

\begin{computationalnote}
В точной арифметике любой ненулевой элемент можно использовать как пивот. При вычислениях с числами с плавающей точкой деление на очень малое число усиливает ошибки округления. Поэтому строки обычно переставляют так, чтобы в текущем столбце использовать элемент наибольшего модуля. Этот приём называется частичным выбором главного элемента.
\end{computationalnote}

\subsection{Совместность и число решений}

После приведения расширенной матрицы к ступенчатому виду возможны три принципиально разные ситуации.

Если появляется строка
\[
\left(
\begin{array}{ccc|c}
0&\cdots&0&c
\end{array}
\right),
\qquad c\neq 0,
\]
то соответствующее уравнение имеет вид \(0=c\). Система несовместна и решений не имеет.

Если противоречивой строки нет и каждый столбец неизвестной содержит пивот, то все неизвестные являются главными. Обратная подстановка даёт единственное решение.

Если противоречивой строки нет, но некоторые неизвестные свободны, то им можно придавать произвольные значения. Главные неизвестные выражаются через свободные, и система имеет семейство решений. Над бесконечным полем, в частности над \(\R\), решений бесконечно много.

\begin{criterion}[критерий чтения ступенчатой системы]
Ступенчатая система несовместна тогда и только тогда, когда содержит уравнение \(0=c\) с \(c\neq 0\). Совместная система имеет единственное решение тогда и только тогда, когда свободных неизвестных нет.
\end{criterion}

\subsection{Однородные системы и геометрия множества решений}

\begin{definition}
Система \(Ax=b\) называется однородной, если \(b=0\), то есть имеет вид
\[
Ax=0.
\]
\end{definition}

Однородная система всегда совместна, поскольку \(x=0\) является решением. Более того, её решения образуют подпространство: если \(Au=0\) и \(Av=0\), то
\[
A(\alpha u+\beta v)
=
\alpha Au+\beta Av
=
0.
\]

\begin{definition}
Множество решений однородной системы называется ядром матрицы \(A\):
\[
\ker A=\{x\in K^n\mid Ax=0\}.
\]
\end{definition}

Для неоднородной совместной системы решения образуют не подпространство, а его сдвиг.

\begin{theorem}
Пусть \(x_0\) --- одно решение системы \(Ax=b\). Тогда множество всех её решений имеет вид
\[
x_0+\ker A
=
\{x_0+z\mid z\in\ker A\}.
\]
\end{theorem}

\begin{proofidea}
Если \(Ax=b\), то
\[
A(x-x_0)=Ax-Ax_0=b-b=0,
\]
следовательно, \(x-x_0\in\ker A\). Обратно, для любого \(z\in\ker A\)
\[
A(x_0+z)=Ax_0+Az=b.
\]
\end{proofidea}

Таким образом, однородная система описывает направления, вдоль которых можно двигаться, не нарушая ограничений. Одно частное решение задаёт положение всего множества решений.

\begin{corollary}
Если однородная система содержит меньше уравнений, чем неизвестных, то после приведения к ступенчатому виду хотя бы одна неизвестная остаётся свободной. Следовательно, система имеет ненулевое решение.
\end{corollary}

Это утверждение предполагает, что речь идёт о конечной системе над полем. Оно не говорит, какое именно ненулевое решение получится: его строят выбором значений свободных переменных.

\subsection{Разобранный пример из курса Винберга}

Рассмотрим систему
\[
\begin{aligned}
x_1+2x_2+x_3&=2,\\
x_1+3x_2+2x_3-x_4&=4,\\
2x_1+x_2-x_3+3x_4&=-2,\\
2x_1-2x_3+3x_4&=1.
\end{aligned}
\]
Её расширенная матрица равна
\[
\left(
\begin{array}{rrrr|r}
1&2&1&0&2\\
1&3&2&-1&4\\
2&1&-1&3&-2\\
2&0&-2&3&1
\end{array}
\right).
\]
Вычтем первую строку из второй, удвоенную первую строку из третьей и четвёртой:
\[
\longrightarrow
\left(
\begin{array}{rrrr|r}
1&2&1&0&2\\
0&1&1&-1&2\\
0&-3&-3&3&-6\\
0&-4&-4&3&-3
\end{array}
\right).
\]
Теперь прибавим к третьей строке утроенную вторую, а к четвёртой --- учетверённую вторую:
\[
\longrightarrow
\left(
\begin{array}{rrrr|r}
1&2&1&0&2\\
0&1&1&-1&2\\
0&0&0&0&0\\
0&0&0&-1&5
\end{array}
\right).
\]
После перестановки двух последних строк получаем ступенчатую форму
\[
\left(
\begin{array}{rrrr|r}
1&2&1&0&2\\
0&1&1&-1&2\\
0&0&0&-1&5\\
0&0&0&0&0
\end{array}
\right).
\]

Пивоты находятся в столбцах \(x_1,x_2,x_4\), поэтому эти переменные главные, а \(x_3\) свободна. Из третьего уравнения
\[
x_4=-5.
\]
Из второго уравнения
\[
x_2+x_3-x_4=2,
\qquad
x_2=-x_3-3.
\]
Из первого уравнения
\[
x_1+2x_2+x_3=2,
\qquad
x_1=x_3+8.
\]
Обозначая свободный параметр через \(t=x_3\), получаем
\[
x=
\begin{pmatrix}
8\\-3\\0\\-5
\end{pmatrix}
+
t
\begin{pmatrix}
1\\-1\\1\\0
\end{pmatrix},
\qquad t\in\R.
\]
Первый столбец --- частное решение неоднородной системы, второй --- базисный вектор ядра. Тем самым вычисление непосредственно реализует формулу
\[
\{x:Ax=b\}=x_0+\ker A.
\]

\paragraph{Практическая интерпретация.}

Пусть неизвестный вектор \(x\) содержит параметры модели, а каждая строка \(A\) задаёт одно линейное измерение. Тогда система \(Ax=b\) отвечает на три разных вопроса. Противоречивая строка показывает, что данные несовместимы с моделью. Единственное решение означает, что измерения однозначно определяют параметры. Свободные переменные показывают направления, которые имеющиеся измерения не различают.

Именно поэтому ядро важно в обратных задачах. Если \(z\in\ker A\), то
\[
A(x+z)=Ax.
\]
Наблюдения не позволяют отличить \(x\) от \(x+z\). Для устранения такой неоднозначности требуются дополнительные измерения, ограничения или критерий выбора решения.

\subsection*{Задачи}

\begin{problem}
\textit{Задача по материалу Винберга.}
Решить систему
\[
\begin{aligned}
x_1+2x_2-x_3&=1,\\
2x_1+4x_2-2x_3&=2,\\
x_1+x_2+x_3&=3
\end{aligned}
\]
и представить множество решений как сумму частного решения и ядра.
\end{problem}

\begin{proof}[Решение]
Расширенная матрица имеет вид
\[
\left(
\begin{array}{rrr|r}
1&2&-1&1\\
2&4&-2&2\\
1&1&1&3
\end{array}
\right).
\]
Вычитаем удвоенную первую строку из второй и первую строку из третьей:
\[
\longrightarrow
\left(
\begin{array}{rrr|r}
1&2&-1&1\\
0&0&0&0\\
0&-1&2&2
\end{array}
\right)
\longrightarrow
\left(
\begin{array}{rrr|r}
1&2&-1&1\\
0&1&-2&-2\\
0&0&0&0
\end{array}
\right).
\]
Переменные \(x_1,x_2\) главные, \(x_3=t\) свободна. Из второй строки
\[
x_2=2t-2,
\]
а из первой
\[
x_1+2(2t-2)-t=1,
\qquad
x_1=5-3t.
\]
Следовательно,
\[
x=
\begin{pmatrix}
5\\-2\\0
\end{pmatrix}
+
t
\begin{pmatrix}
-3\\2\\1
\end{pmatrix}.
\]
Вектор \((5,-2,0)^\mathsf{T}\) является частным решением, а
\[
\ker A
=
\operatorname{span}
\left\{
\begin{pmatrix}
-3\\2\\1
\end{pmatrix}
\right\}.
\]
Вторая строка исходной системы не добавляла информации: она была удвоенной первой строкой.
\end{proof}

\begin{problem}
\textit{Практическая задача: проверка совместности.}
Определить, имеет ли решение система
\[
\begin{aligned}
x_1+x_2+x_3&=1,\\
2x_1+2x_2+2x_3&=3.
\end{aligned}
\]
\end{problem}

\begin{proof}[Решение]
Вычтем из второй строки удвоенную первую:
\[
\left(
\begin{array}{rrr|r}
1&1&1&1\\
2&2&2&3
\end{array}
\right)
\longrightarrow
\left(
\begin{array}{rrr|r}
1&1&1&1\\
0&0&0&1
\end{array}
\right).
\]
Вторая строка означает \(0=1\), поэтому система несовместна. Геометрически левые части задают одну и ту же плоскость, но правые части требуют двух разных значений одного и того же линейного измерения.
\end{proof}

\begin{problem}
\textit{Практическая задача: пространство ненаблюдаемых направлений.}
Линейное измерение задано матрицей
\[
A=
\begin{pmatrix}
1&1&0&0\\
0&1&1&0\\
0&0&1&1
\end{pmatrix}.
\]
Найти все изменения параметров \(z\in\R^4\), которые не меняют результат измерения, то есть удовлетворяют \(Az=0\).
\end{problem}

\begin{proof}[Решение]
Система имеет вид
\[
\begin{aligned}
z_1+z_2&=0,\\
z_2+z_3&=0,\\
z_3+z_4&=0.
\end{aligned}
\]
Положим \(z_4=t\). Тогда
\[
z_3=-t,
\qquad
z_2=t,
\qquad
z_1=-t.
\]
Поэтому
\[
z=t
\begin{pmatrix}
-1\\1\\-1\\1
\end{pmatrix},
\qquad
\ker A=
\operatorname{span}
\left\{
\begin{pmatrix}
-1\\1\\-1\\1
\end{pmatrix}
\right\}.
\]
Если параметр модели изменить вдоль этого направления, наблюдаемый вектор \(Ax\) останется прежним. Это точная линейная форма неидентифицируемости.
\end{proof}

\begin{problem}
\textit{Практическая задача: восстановление параметров.}
Аффинная функция
\[
f(t)=\alpha t+\beta
\]
проходит через точки \((1,3)\) и \((4,9)\). Записать задачу как линейную систему и найти параметры.
\end{problem}

\begin{proof}[Решение]
Подстановка двух наблюдений даёт
\[
\begin{pmatrix}
1&1\\
4&1
\end{pmatrix}
\begin{pmatrix}
\alpha\\\beta
\end{pmatrix}
=
\begin{pmatrix}
3\\9
\end{pmatrix}.
\]
Вычитая первую строку из второй, получаем
\[
3\alpha=6,
\qquad
\alpha=2.
\]
Тогда из первого уравнения
\[
\beta=1.
\]
Следовательно,
\[
f(t)=2t+1.
\]
Два независимых точных измерения однозначно определили два параметра. При большем числе несовместимых наблюдений точного решения может не существовать; тогда возникает задача наименьших квадратов, которая будет разобрана после введения ортогональных проекций.
\end{proof}

\section{Ранг матрицы, ядро и образ}
\label{sec:rank-kernel-image}

Ранг показывает, сколько независимых направлений действительно передаёт матрица. Размер матрицы сообщает лишь число входных и выходных координат; ранг отделяет содержательную размерность преобразования от возможной избыточности строк и столбцов.

\subsection{Строчное и столбцовое пространства}

Пусть \(A\in K^{m\times n}\). Её строки являются векторами пространства \(K^n\), а столбцы --- векторами пространства \(K^m\).

\begin{definition}
Строчным пространством матрицы \(A\) называется линейная оболочка её строк. Столбцовым пространством называется линейная оболочка её столбцов:
\[
\operatorname{Row}(A)=\operatorname{span}\{\text{строки }A\},
\qquad
\operatorname{Col}(A)=\operatorname{span}\{\text{столбцы }A\}.
\]
\end{definition}

Винберг сначала определяет ранг как размерность строчного пространства, а затем доказывает, что размерность столбцового пространства совпадает с ней.

\begin{theorem}
Для любой матрицы \(A\)
\[
\dim \operatorname{Row}(A)
=
\dim \operatorname{Col}(A).
\]
Общее значение называется рангом матрицы и обозначается
\[
\operatorname{rank}A.
\]
\end{theorem}

Равенство важно не только теоретически. Строки описывают независимые линейные ограничения или измерения, а столбцы --- независимые направления результата, возникающие при изменении входных координат. Одна и та же величина измеряет обе стороны матрицы.

Элементарные преобразования строк не меняют ранг. Поэтому ранг равен числу ненулевых строк ступенчатой формы. Если ведущие элементы ступенчатой формы находятся в столбцах с номерами \(j_1,\dots,j_r\), то столбцы исходной матрицы с этими номерами образуют базис \(\operatorname{Col}(A)\). Здесь необходимо брать столбцы именно исходной матрицы: преобразования строк меняют сами столбцы, хотя сохраняют зависимости между ними.

\subsection{Матрица как линейное отображение}

Матрица \(A\in K^{m\times n}\) задаёт отображение
\[
T_A\colon K^n\to K^m,
\qquad
T_A(x)=Ax.
\]
Оно линейно, поскольку
\[
A(\alpha x+\beta y)=\alpha Ax+\beta Ay.
\]

\begin{definition}
Образом матрицы \(A\) называется множество всех возможных результатов:
\[
\operatorname{im}A
=
\{Ax\mid x\in K^n\}.
\]
Ядром называется множество входов, переходящих в нуль:
\[
\ker A
=
\{x\in K^n\mid Ax=0\}.
\]
\end{definition}

Если \(a_1,\dots,a_n\) --- столбцы \(A\), то
\[
Ax=x_1a_1+\cdots+x_na_n.
\]
Следовательно,
\[
\operatorname{im}A=\operatorname{Col}(A),
\qquad
\dim\operatorname{im}A=\operatorname{rank}A.
\]
Образ содержит все достижимые выходы, а ядро --- все изменения входа, которые преобразование не замечает.

\begin{theorem}[теорема о ранге и дефекте]
Для матрицы \(A\in K^{m\times n}\)
\[
\dim\ker A+\operatorname{rank}A=n.
\]
\end{theorem}

\begin{proofidea}
После приведения \(A\) к ступенчатому виду число главных переменных равно числу пивотов, то есть рангу \(A\). Остальные \(n-\operatorname{rank}A\) переменных свободны и дают столько же базисных направлений ядра.
\end{proofidea}

Формула разделяет \(n\) входных степеней свободы на две группы. Ранг считает направления, которые проявляются в результате, а размерность ядра --- направления, теряемые преобразованием. Это не утверждение о буквальном разбиении координат исходного базиса: соответствующие направления обычно являются их линейными комбинациями.

\subsection{Совместность, единственность и обратимость}

Система \(Ax=b\) имеет решение тогда и только тогда, когда \(b\in\operatorname{im}A\). В координатной форме это даёт критерий Кронекера--Капелли:
\[
Ax=b \text{ совместна}
\quad\Longleftrightarrow\quad
\operatorname{rank}A
=
\operatorname{rank}(A\mid b).
\]
Если система совместна, то решение единственно тогда и только тогда, когда \(\ker A=\{0\}\), то есть
\[
\operatorname{rank}A=n.
\]
Для квадратной матрицы порядка \(n\) это равносильно обратимости. Тогда каждый выход \(b\in K^n\) достигается единственным входом
\[
x=A^{-1}b.
\]

При композиции преобразований информация не может самопроизвольно увеличиться. Для согласованных матриц \(A\) и \(B\)
\[
\operatorname{rank}(AB)
\leq
\min\{\operatorname{rank}A,\operatorname{rank}B\}.
\]
Матрица \(B\) сначала ограничивает множество промежуточных состояний своим образом, а матрица \(A\) затем может дополнительно склеить некоторые из этих состояний.

\subsection{Одно вычисление: ранг, образ и ядро}

Рассмотрим матрицу
\[
A=
\begin{pmatrix}
1&0&1&2\\
0&1&1&1\\
1&1&2&3
\end{pmatrix}.
\]
Третья строка равна сумме первых двух. После замены
\(R_3\leftarrow R_3-R_1-R_2\) получаем
\[
\begin{pmatrix}
1&0&1&2\\
0&1&1&1\\
0&0&0&0
\end{pmatrix}.
\]
Отсюда
\[
\operatorname{rank}A=2.
\]
Пивоты стоят в первом и втором столбцах, поэтому базис образа составляют соответствующие столбцы исходной матрицы:
\[
\operatorname{im}A
=
\operatorname{span}
\left\{
\begin{pmatrix}1\\0\\1\end{pmatrix},
\begin{pmatrix}0\\1\\1\end{pmatrix}
\right\}.
\]

Для ядра решаем \(Ax=0\):
\[
x_1+x_3+2x_4=0,
\qquad
x_2+x_3+x_4=0.
\]
Положив \(x_3=s\), \(x_4=t\), получаем
\[
x=
 s
\begin{pmatrix}-1\\-1\\1\\0\end{pmatrix}
+t
\begin{pmatrix}-2\\-1\\0\\1\end{pmatrix}.
\]
Следовательно,
\[
\dim\ker A=2,
\qquad
\dim\ker A+\operatorname{rank}A=2+2=4.
\]
Одно приведение к ступенчатому виду одновременно выявило число независимых выходных направлений и все невидимые изменения входа.

\subsection*{Задачи}

\begin{problem}[Задача Винберга]
Доказать, что для матриц одинакового размера
\[
\operatorname{rank}(A+B)
\leq
\operatorname{rank}A+
\operatorname{rank}B.
\]
Привести пример, в котором достигается равенство.
\end{problem}

\begin{proof}[Решение]
Каждый столбец \(A+B\) является суммой соответствующего столбца \(A\) и столбца \(B\). Поэтому
\[
\operatorname{Col}(A+B)
\subseteq
\operatorname{Col}(A)+\operatorname{Col}(B).
\]
Размерность суммы двух подпространств не превосходит суммы их размерностей, откуда
\[
\operatorname{rank}(A+B)
\leq
\operatorname{rank}A+
\operatorname{rank}B.
\]
Равенство достигается, например, для
\[
A=
\begin{pmatrix}1&0\\0&0\end{pmatrix},
\qquad
B=
\begin{pmatrix}0&0\\0&1\end{pmatrix}.
\]
Обе матрицы имеют ранг \(1\), а \(A+B=I_2\) имеет ранг \(2\).
\end{proof}

\begin{problem}[Задача по материалу Винберга]
Для согласованных матриц \(A\) и \(B\) доказать
\[
\operatorname{rank}(AB)
\leq
\min\{\operatorname{rank}A,\operatorname{rank}B\}.
\]
\end{problem}

\begin{proof}[Решение]
Любой результат композиции имеет вид
\[
ABx=A(Bx),
\]
поэтому \(\operatorname{im}(AB)\subseteq\operatorname{im}A\). Значит,
\[
\operatorname{rank}(AB)\leq\operatorname{rank}A.
\]
С другой стороны, отображение \(A\) действует лишь на векторы из \(\operatorname{im}B\). Линейное отображение не может увеличить размерность пространства, следовательно,
\[
\dim A(\operatorname{im}B)
\leq
\dim\operatorname{im}B.
\]
Но \(A(\operatorname{im}B)=\operatorname{im}(AB)\), поэтому
\[
\operatorname{rank}(AB)
\leq
\operatorname{rank}B.
\]
Совмещая две оценки, получаем требуемое неравенство.
\end{proof}

\begin{problem}[Практическая задача: зависимые датчики]
Изменение состояния помещения описывается вектором
\[
x=(T,H,C)^{\trans},
\]
где \(T\) --- изменение температуры, \(H\) --- влажности, \(C\) --- концентрации углекислого газа. Три датчика выдают
\[
y=
\begin{pmatrix}
1&1&0\\
0&1&1\\
1&2&1
\end{pmatrix}x.
\]
Сколько независимых измерений дают датчики и какие изменения состояния они не способны различить?
\end{problem}

\begin{proof}[Решение]
Третья строка матрицы равна сумме первых двух, поэтому третий датчик не добавляет независимой информации:
\[
\operatorname{rank}A=2.
\]
Невидимые изменения находятся из \(Ax=0\):
\[
T+H=0,
\qquad
H+C=0.
\]
Полагая \(H=t\), получаем
\[
x=t
\begin{pmatrix}-1\\1\\-1\end{pmatrix}.
\]
Таким образом, одновременное изменение температуры, влажности и концентрации в пропорции
\[
(-1,1,-1)
\]
не меняет ни одного показания. Система имеет три физических параметра, но наблюдает только две независимые их комбинации. Для полного восстановления состояния требуется датчик с новым, линейно независимым законом измерения.
\end{proof}

\begin{problem}[Практическая задача: достижимая деформация]
Два привода управляют вертикальными смещениями трёх контрольных точек гибкой панели. При командах \(u_1,u_2\) смещения равны
\[
d=
\begin{pmatrix}
1&0\\
0&1\\
1&1
\end{pmatrix}
\begin{pmatrix}u_1\\u_2\end{pmatrix}.
\]
Проверить достижимость профилей
\[
d^{(1)}=(2,-1,1)^{\trans},
\qquad
 d^{(2)}=(2,-1,0)^{\trans}.
\]
Для достижимого профиля найти команды приводов.
\end{problem}

\begin{proof}[Решение]
Столбцы матрицы линейно независимы, поэтому её ранг равен \(2\), а ядро нулевое. Любой достижимый профиль обязан удовлетворять связи
\[
d_3=d_1+d_2.
\]
Для первого профиля
\[
1=2+(-1),
\]
поэтому он достижим. Первые две строки сразу дают
\[
u_1=2,
\qquad
u_2=-1.
\]
Для второго профиля требовалось бы \(0=2+(-1)=1\), что невозможно. Его нельзя получить никакими командами: ограничение связано не с выбором алгоритма, а с двумерным образом матрицы приводов.
\end{proof}

\section{Определитель, объём и обратимость}
\label{sec:determinant}

Определитель сопоставляет квадратной матрице одно число. Это число одновременно отвечает на два вопроса: сохраняет ли линейное преобразование все направления и во сколько раз оно изменяет ориентированный объём. Поэтому определитель служит компактным критерием вырожденности, но не заменяет саму матрицу и не описывает преобразование полностью.

\subsection{Геометрический смысл}

Пусть столбцы матрицы
\[
A=\begin{pmatrix}a&c\\ b&d\end{pmatrix}
\]
задают векторы
\[
v_1=(a,b)^{\trans},
\qquad
v_2=(c,d)^{\trans}.
\]
Ориентированная площадь параллелограмма, натянутого на эти векторы, равна
\[
\det A=ad-bc.
\]
Абсолютная величина \(\lvert\det A\rvert\) является обычной площадью. Знак показывает ориентацию: положительный знак означает сохранение принятой ориентации, отрицательный --- её обращение. При \(\det A=0\) площадь исчезает, поскольку столбцы коллинеарны.

В пространстве \(\R^3\) величина \(\lvert\det A\rvert\) равна объёму параллелепипеда, натянутого на столбцы \(A\). В размерности \(n\) сохраняется тот же смысл: определитель измеряет ориентированный \(n\)-мерный объём.

\begin{application}
Если линейное преобразование задаётся формулой \(y=Ax\), то небольшой объём в пространстве входов после преобразования умножается на \(\lvert\det A\rvert\). При нулевом определителе весь \(n\)-мерный объём сплющивается в множество меньшей размерности, поэтому восстановить вход по выходу однозначно невозможно.
\end{application}

\subsection{Строгое определение}

Геометрический образ площади и объёма распространяется на произвольное поле через три свойства.

\begin{definition}
Определителем квадратной матрицы порядка \(n\) называется единственная функция её столбцов
\[
\det(v_1,\dots,v_n),
\]
которая линейна по каждому столбцу, меняет знак при перестановке двух столбцов и удовлетворяет нормировке
\[
\det I_n=1.
\]
\end{definition}

Линейность по каждому аргументу означает, например,
\[
\det(v_1,\dots,\alpha u+\beta w,\dots,v_n)
=
\alpha\det(v_1,\dots,u,\dots,v_n)
+
\beta\det(v_1,\dots,w,\dots,v_n).
\]
Из кососимметричности следует, что при совпадении двух столбцов определитель равен нулю.

Для матрицы \(A=(a_{ij})\) это определение эквивалентно формуле Лейбница
\[
\det A
=
\sum_{\sigma\in S_n}
\operatorname{sgn}(\sigma)
\prod_{i=1}^{n}a_{i,\sigma(i)},
\]
где \(S_n\) --- множество перестановок чисел \(1,\dots,n\), а \(\operatorname{sgn}(\sigma)\) равен \(1\) или \(-1\) в зависимости от чётности перестановки. Формула важна как точное описание, но для вычислений при больших \(n\) непригодна: она содержит \(n!\) слагаемых.

\subsection{Свойства, необходимые в вычислениях}

Перестановка двух строк или столбцов меняет знак определителя. Умножение одной строки на число \(\lambda\) умножает определитель на \(\lambda\). Прибавление к строке другой строки, умноженной на число, определитель не меняет. Поэтому определитель удобно вычислять теми же преобразованиями, что используются в методе Гаусса.

\begin{proposition}
Для треугольной матрицы
\[
T=\begin{pmatrix}
 t_{11}&*&\cdots&*\\
 0&t_{22}&\cdots&*\\
 \vdots&\ddots&\ddots&\vdots\\
 0&\cdots&0&t_{nn}
\end{pmatrix}
\]
выполнено
\[
\det T=t_{11}t_{22}\cdots t_{nn}.
\]
\end{proposition}

Таким образом, для практического вычисления матрицу приводят к треугольному виду, отслеживая перестановки и умножения строк. Именно этот принцип лежит в основе вычисления определителя через разложение \(LU\).

\begin{theorem}
Для квадратных матриц одинакового порядка выполняются равенства
\[
\det(A^{\trans})=\det A,
\qquad
\det(AB)=\det A\,\det B.
\]
\end{theorem}

Второе равенство выражает последовательное изменение объёма: сначала преобразование \(B\) умножает объём на \(\lvert\det B\rvert\), затем преобразование \(A\) --- на \(\lvert\det A\rvert\).

\subsection{Определитель и обратимость}

\begin{theorem}[критерий невырожденности]
Для квадратной матрицы \(A\in K^{n\times n}\) следующие условия равносильны:
\[
\det A\neq 0,
\qquad
\operatorname{rank}A=n,
\qquad
\ker A=\{0\},
\qquad
A^{-1}\text{ существует}.
\]
\end{theorem}

\begin{proofidea}
После приведения матрицы к ступенчатому виду она имеет полный ранг тогда и только тогда, когда все диагональные элементы полученной треугольной матрицы ненулевые. Это равносильно ненулевому произведению диагональных элементов, то есть ненулевому определителю. Остальные эквивалентности уже следуют из связи полного ранга, нулевого ядра и обратимости.
\end{proofidea}

Если \(A\) обратима, то из мультипликативности следует
\[
1=\det I=\det(AA^{-1})=\det A\,\det A^{-1},
\]
поэтому
\[
\det A^{-1}=\frac{1}{\det A}.
\]

\subsection{Миноры, алгебраические дополнения и формулы Крамера}

Пусть \(A=(a_{ij})\) --- квадратная матрица. Удалив строку \(i\) и столбец \(j\), получаем дополнительный минор \(M_{ij}\). Алгебраическим дополнением элемента \(a_{ij}\) называется
\[
C_{ij}=(-1)^{i+j}M_{ij}.
\]
Определитель можно разложить по любой строке или столбцу:
\[
\det A=\sum_{j=1}^{n}a_{ij}C_{ij},
\qquad
\det A=\sum_{i=1}^{n}a_{ij}C_{ij}.
\]
Разложение удобно для небольших разреженных матриц, но не является основным численным методом для плотных матриц большого порядка.

Для системы
\[
Ax=b,
\qquad
\det A\neq 0,
\]
обозначим через \(A_i\) матрицу, в которой \(i\)-й столбец \(A\) заменён столбцом \(b\). Тогда формулы Крамера имеют вид
\[
x_i=\frac{\det A_i}{\det A}.
\]
Они дают точную теоретическую формулу и полезны для систем порядка \(2\) или \(3\), но в вычислительной практике систему решают методом Гаусса или факторизациями матриц.

Для матрицы второго порядка получается знакомая формула
\[
\begin{pmatrix}a&b\\c&d\end{pmatrix}^{-1}
=
\frac{1}{ad-bc}
\begin{pmatrix}d&-b\\-c&a\end{pmatrix},
\qquad
ad-bc\neq 0.
\]

\subsection*{Задачи}

\begin{problem}[Задача Винберга: несовместная система]
Пусть \(A_i\) получается из квадратной матрицы \(A\) заменой \(i\)-го столбца столбцом \(b\). Доказать: если
\[
\det A=0,
\qquad
\det A_i\neq 0
\]
хотя бы для одного \(i\), то система \(Ax=b\) несовместна.
\end{problem}

\begin{proof}[Решение]
Предположим, что решение \(x=(x_1,\dots,x_n)^{\trans}\) существует. Если \(a_1,\dots,a_n\) --- столбцы \(A\), то
\[
b=x_1a_1+\cdots+x_na_n.
\]
Подставим это выражение вместо \(i\)-го столбца и воспользуемся линейностью определителя по столбцу. Все слагаемые, кроме содержащего \(x_i a_i\), имеют два одинаковых столбца и потому равны нулю. Следовательно,
\[
\det A_i=x_i\det A=0,
\]
что противоречит условию \(\det A_i\neq 0\). Значит, решения нет.
\end{proof}

\begin{problem}[Задача Винберга: целочисленная обратная матрица]
Пусть \(A\) --- невырожденная квадратная матрица с целыми элементами. Доказать, что \(A^{-1}\) имеет целые элементы тогда и только тогда, когда
\[
\det A=\pm 1.
\]
\end{problem}

\begin{proof}[Решение]
Если \(A^{-1}\) целочисленна, то \(\det A\) и \(\det A^{-1}\) --- целые числа. Из
\[
\det A\,\det A^{-1}=1
\]
следует \(\det A=\pm1\).

Обратно, элементы присоединённой матрицы являются алгебраическими дополнениями и потому целые. Формула
\[
A^{-1}=\frac{1}{\det A}\operatorname{adj}A
\]
при \(\det A=\pm1\) не создаёт дробных элементов. Следовательно, \(A^{-1}\) целочисленна.
\end{proof}

\begin{problem}[Практическая задача: контроль качества треугольного элемента сетки]
В двумерной расчётной сетке треугольник задан вершинами
\[
p_0=(0,0),
\qquad
p_1=(3,1),
\qquad
p_2=(1,4).
\]
Найти его площадь и ориентацию. Затем проверить, что произойдёт, если третья вершина переместится в точку \(p'_2=(6,2)\).
\end{problem}

\begin{proof}[Решение]
Матрица рёбер исходного треугольника равна
\[
E=
\begin{pmatrix}
3&1\\
1&4
\end{pmatrix}.
\]
Её определитель
\[
\det E=3\cdot4-1\cdot1=11>0.
\]
Площадь треугольника составляет половину площади параллелограмма:
\[
S=\frac12\lvert\det E\rvert=\frac{11}{2}.
\]
Положительный знак показывает принятую положительную ориентацию вершин.

После перемещения вершины
\[
E'=
\begin{pmatrix}
3&6\\
1&2
\end{pmatrix},
\qquad
\det E'=6-6=0.
\]
Треугольник вырождается в отрезок. Такой элемент нельзя использовать в конечно-элементной сетке: локальное преобразование координат теряет обратимость.
\end{proof}

\begin{problem}[Практическая задача: два привода плоского позиционера]
Команды двух приводов \(u_1,u_2\) создают перемещение рабочего наконечника
\[
d=
\begin{pmatrix}
2&1\\
1&2
\end{pmatrix}
\begin{pmatrix}u_1\\u_2\end{pmatrix}
\quad\text{миллиметров}.
\]
Проверить, можно ли независимо управлять двумя координатами, и найти команды для перемещения
\[
d=(5,4)^{\trans}.
\]
\end{problem}

\begin{proof}[Решение]
Матрица управления имеет определитель
\[
\det A=2\cdot2-1\cdot1=3\neq0.
\]
Следовательно, два режима приводов линейно независимы, а любое достаточно малое плоское перемещение имеет единственный набор команд. Обратная матрица равна
\[
A^{-1}=\frac13
\begin{pmatrix}
2&-1\\
-1&2
\end{pmatrix}.
\]
Поэтому
\[
\begin{pmatrix}u_1\\u_2\end{pmatrix}
=A^{-1}d
=
\frac13
\begin{pmatrix}
2&-1\\
-1&2
\end{pmatrix}
\begin{pmatrix}5\\4\end{pmatrix}
=
\begin{pmatrix}2\\1\end{pmatrix}.
\]
Команды \(u_1=2\) и \(u_2=1\) дают требуемое перемещение. Величина \(\lvert\det A\rvert=3\) также означает, что единичный квадрат в пространстве команд переходит в параллелограмм площади \(3\) в пространстве перемещений.
\end{proof}

\section{Линейные отображения и их матрицы}
\label{sec:linear-maps}

Матрица является координатной записью линейного правила. Само правило существует независимо от выбора базисов: оно переводит векторы одного пространства в векторы другого и сохраняет линейные комбинации. Разделение отображения и его матрицы важно, поскольку одно и то же преобразование в разных координатах записывается разными таблицами чисел.

\subsection{Определение линейного отображения}

\begin{definition}
Пусть \(V\) и \(U\) --- векторные пространства над одним полем \(K\).
Отображение
\[
\varphi\colon V\to U
\]
называется линейным, если для любых \(x,y\in V\) и
\(\alpha,\beta\in K\)
\[
\varphi(\alpha x+\beta y)
=
\alpha\varphi(x)+\beta\varphi(y).
\]
\end{definition}

Из определения сразу следуют равенства
\[
\varphi(0)=0,
\qquad
\varphi(-x)=-\varphi(x).
\]
Поэтому параллельный перенос не является линейным отображением, если он сдвигает начало координат. Такие преобразования будут рассматриваться позже как аффинные.

Типичными линейными отображениями являются умножение на матрицу, ортогональное проектирование, дифференцирование функций и переход от коэффициентов модели к её наблюдаемому отклику.

\subsection{Отображение определяется образами базиса}

Пусть \(e_1,\dots,e_n\) --- базис пространства \(V\). Каждый вектор имеет единственное разложение
\[
x=x_1e_1+\cdots+x_ne_n.
\]
Линейность даёт
\[
\varphi(x)
=
x_1\varphi(e_1)+\cdots+x_n\varphi(e_n).
\]

\begin{proposition}
Линейное отображение однозначно определяется образами базисных векторов. Более того, произвольный набор векторов
\[
u_1,\dots,u_n\in U
\]
задаёт единственное линейное отображение, для которого
\[
\varphi(e_j)=u_j.
\]
\end{proposition}

Это утверждение задаёт практический способ построения преобразования: достаточно указать, что происходит с независимыми базисными направлениями. Поведение на всех остальных векторах восстанавливается линейностью.

\subsection{Матрица линейного отображения}

Выберем базис
\[
e=(e_1,\dots,e_n)
\]
в пространстве \(V\) и базис
\[
f=(f_1,\dots,f_m)
\]
в пространстве \(U\). Разложим образы базисных векторов:
\[
\varphi(e_j)
=
\sum_{i=1}^{m}a_{ij}f_i.
\]

\begin{definition}
Матрицей линейного отображения \(\varphi\colon V\to U\) в базисах
\(e\) и \(f\) называется матрица
\[
A=(a_{ij})\in K^{m\times n}.
\]
Её \(j\)-й столбец является столбцом координат вектора
\(\varphi(e_j)\) в базисе \(f\).
\end{definition}

Для любого \(x\in V\) выполняется основная координатная формула
\[
[\varphi(x)]_f
=
A[x]_e.
\]
Размер матрицы определяется направлениями отображения:
число столбцов равно \(\dim V\), а число строк равно \(\dim U\).

\begin{example}
Пусть \(P_2\) --- пространство многочленов степени не выше двух,
а
\[
\varphi\colon P_2\to\R^2,
\qquad
\varphi(p)=\bigl(p(0),p(1)\bigr).
\]
В базисе \(e=(1,t,t^2)\) пространства \(P_2\) и стандартном базисе
\(\R^2\)
\[
\varphi(1)=(1,1),
\qquad
\varphi(t)=(0,1),
\qquad
\varphi(t^2)=(0,1).
\]
Следовательно,
\[
A=
\begin{pmatrix}
1&0&0\\
1&1&1
\end{pmatrix}.
\]
Для многочлена \(p(t)=2-t+3t^2\)
\[
[p]_e=
\begin{pmatrix}
2\\-1\\3
\end{pmatrix},
\qquad
A[p]_e
=
\begin{pmatrix}
2\\4
\end{pmatrix}.
\]
Это совпадает с прямым вычислением \(p(0)=2\), \(p(1)=4\).
\end{example}

\subsection{Ядро, образ и изоморфизм}

Для абстрактного линейного отображения определения имеют тот же смысл, что и для матрицы:
\[
\ker\varphi
=
\{x\in V\mid \varphi(x)=0\},
\qquad
\operatorname{im}\varphi
=
\{\varphi(x)\mid x\in V\}.
\]
Оба множества являются подпространствами. В выбранных базисах
\[
\dim\operatorname{im}\varphi=\operatorname{rank}A,
\]
а для конечномерного \(V\)
\[
\dim\ker\varphi+\dim\operatorname{im}\varphi=\dim V.
\]

Отображение инъективно тогда и только тогда, когда
\[
\ker\varphi=\{0\}.
\]
Оно сюръективно тогда и только тогда, когда
\[
\operatorname{im}\varphi=U.
\]

\begin{definition}
Биективное линейное отображение называется изоморфизмом векторных пространств.
\end{definition}

Изоморфизм сохраняет линейную структуру и допускает обратное линейное отображение. В конечномерных пространствах одинаковой размерности достаточно проверить одно из двух свойств: инъективность или сюръективность.

\subsection{Композиция отображений}

Пусть
\[
\psi\colon W\to V,
\qquad
\varphi\colon V\to U.
\]
Их композиция
\[
\varphi\circ\psi\colon W\to U
\]
сначала применяет \(\psi\), затем \(\varphi\). Композиция линейных отображений линейна.

Если в согласованных базисах матрица \(\psi\) равна \(B\), а матрица
\(\varphi\) равна \(A\), то
\[
[\varphi\circ\psi]=AB.
\]
Порядок множителей соответствует порядку применения справа налево:
\[
x
\longmapsto
Bx
\longmapsto
A(Bx).
\]
Вообще \(AB\neq BA\), поэтому перестановка этапов обработки обычно меняет результат.

\subsection{Смена базисов}

Пусть \(A\) --- матрица отображения \(\varphi\colon V\to U\) в старых базисах \(e\) и \(f\). Пусть \(C\) переводит новые координаты входа в старые:
\[
[x]_e=C[x]_{e'},
\]
а \(D\) переводит новые координаты выхода в старые:
\[
[y]_f=D[y]_{f'}.
\]
Тогда матрица того же отображения в новых базисах равна
\[
A'=D^{-1}AC.
\]

Правая матрица \(C\) сначала переводит вход из новых координат в старые. Матрица \(A\) выполняет само преобразование. Левая матрица \(D^{-1}\) переводит полученный выход в новый базис. Меняется координатное описание, но не действие отображения на векторы.

Если \(V=U\) и на входе и выходе меняется один и тот же базис, то
\[
A'=C^{-1}AC.
\]
Такие матрицы называются подобными. Позже выбор подходящего базиса позволит приводить операторы к диагональной или жордановой форме.

\subsection*{Задачи}

\begin{problem}[Задача по материалу Винберга: дифференцирование]
Пусть \(P_3\) и \(P_2\) --- пространства вещественных многочленов степени не выше \(3\) и \(2\). Для отображения
\[
D\colon P_3\to P_2,
\qquad
D(p)=p'
\]
найти матрицу в стандартных базисах, а также ядро, образ и ранг.
\end{problem}

\begin{proof}[Решение]
Используем базисы
\[
(1,t,t^2,t^3)
\quad\text{и}\quad
(1,t,t^2).
\]
Образы базисных векторов равны
\[
D(1)=0,
\qquad
D(t)=1,
\qquad
D(t^2)=2t,
\qquad
D(t^3)=3t^2.
\]
Их координаты образуют столбцы матрицы:
\[
[D]=
\begin{pmatrix}
0&1&0&0\\
0&0&2&0\\
0&0&0&3
\end{pmatrix}.
\]
Производная равна нулю только у постоянных многочленов, поэтому
\[
\ker D=\operatorname{span}\{1\}.
\]
Каждый многочлен
\[
a+bt+ct^2
\]
является производной многочлена
\[
at+\frac{b}{2}t^2+\frac{c}{3}t^3,
\]
следовательно,
\[
\operatorname{im}D=P_2,
\qquad
\operatorname{rank}D=3.
\]
Равенство
\[
\dim\ker D+\dim\operatorname{im}D=1+3=4
\]
совпадает с \(\dim P_3\).
\end{proof}

\begin{problem}[Задача по материалу Винберга: композиция отображений]
Линейные отображения
\[
\varphi\colon\R^2\to\R^3,
\qquad
\psi\colon\R^3\to\R^2
\]
имеют в стандартных базисах матрицы
\[
A=
\begin{pmatrix}
1&0\\
0&1\\
1&1
\end{pmatrix},
\qquad
B=
\begin{pmatrix}
1&0&1\\
0&1&1
\end{pmatrix}.
\]
Найти матрицу \(\psi\circ\varphi\) и проверить результат на векторе
\(x=(1,2)^{\trans}\).
\end{problem}

\begin{proof}[Решение]
Сначала применяется \(\varphi\), затем \(\psi\), поэтому матрица композиции равна
\[
BA=
\begin{pmatrix}
1&0&1\\
0&1&1
\end{pmatrix}
\begin{pmatrix}
1&0\\
0&1\\
1&1
\end{pmatrix}
=
\begin{pmatrix}
2&1\\
1&2
\end{pmatrix}.
\]
Отсюда
\[
(\psi\circ\varphi)(x)
=
\begin{pmatrix}
2&1\\
1&2
\end{pmatrix}
\begin{pmatrix}
1\\2
\end{pmatrix}
=
\begin{pmatrix}
4\\5
\end{pmatrix}.
\]
Проверим по этапам:
\[
\varphi(x)
=
A
\begin{pmatrix}
1\\2
\end{pmatrix}
=
\begin{pmatrix}
1\\2\\3
\end{pmatrix},
\qquad
\psi(\varphi(x))
=
B
\begin{pmatrix}
1\\2\\3
\end{pmatrix}
=
\begin{pmatrix}
4\\5
\end{pmatrix}.
\]
Результаты совпали.
\end{proof}

\begin{problem}[Практическая задача: яркость цифрового цвета]
Яркость RGB-цвета приближённо вычисляется линейным отображением
\[
L(R,G,B)
=
\frac{3}{10}R+\frac{6}{10}G+\frac{1}{10}B.
\]
Найти яркость цвета \((120,80,40)\) и ненулевое изменение цвета, которое не меняет яркость.
\end{problem}

\begin{proof}[Решение]
Для заданного цвета
\[
L(120,80,40)
=
36+48+4
=
88.
\]
Изменение \((\Delta R,\Delta G,\Delta B)\) невидимо для яркостного канала, если оно лежит в ядре:
\[
3\Delta R+6\Delta G+\Delta B=0.
\]
Например,
\[
(\Delta R,\Delta G,\Delta B)=(2,-1,0)
\]
удовлетворяет этому условию, поскольку
\[
3\cdot2+6\cdot(-1)+0=0.
\]
Следовательно, цвета, отличающиеся на \((2,-1,0)\), имеют одинаковую яркость в этой модели, хотя их оттенки различаются. Одномерный выход неизбежно теряет часть информации о трёхмерном цвете.
\end{proof}

\begin{problem}[Практическая задача: порядок геометрических преобразований]
Точка плоской детали сначала растягивается по оси \(x\) в \(2\) раза и сжимается по оси \(y\) в \(2\) раза, затем поворачивается на \(90^\circ\) против часовой стрелки. Найти итоговую матрицу и образ точки
\[
p=(1,2)^{\trans}.
\]
Сравнить с результатом, полученным при обратном порядке операций.
\end{problem}

\begin{proof}[Решение]
Матрицы масштабирования и поворота равны
\[
S=
\begin{pmatrix}
2&0\\
0&\frac12
\end{pmatrix},
\qquad
R=
\begin{pmatrix}
0&-1\\
1&0
\end{pmatrix}.
\]
Сначала применяется \(S\), затем \(R\), поэтому итоговая матрица
\[
T=RS
=
\begin{pmatrix}
0&-\frac12\\
2&0
\end{pmatrix}.
\]
Для точки \(p\)
\[
Tp
=
\begin{pmatrix}
0&-\frac12\\
2&0
\end{pmatrix}
\begin{pmatrix}
1\\2
\end{pmatrix}
=
\begin{pmatrix}
-1\\2
\end{pmatrix}.
\]

При обратном порядке
\[
SR
=
\begin{pmatrix}
0&-2\\
\frac12&0
\end{pmatrix},
\qquad
SRp=
\begin{pmatrix}
-4\\
\frac12
\end{pmatrix}.
\]
Результаты различаются. В геометрическом конвейере порядок матриц кодирует порядок реальных операций над формой и не может быть произвольно переставлен.
\end{proof}

\section{Линейные функционалы и сопряжённое пространство}
\label{sec:dual-space}

До сих пор линейные отображения переводили векторы в другие векторы. Особый случай возникает, когда результатом является одно число. Такое отображение измеряет некоторую линейную характеристику состояния: суммарную нагрузку, яркость, поток, работу силы или первую вариацию целевой функции.

\subsection{Линейный функционал}

\begin{definition}
Линейным функционалом на векторном пространстве \(V\) над полем \(K\) называется линейное отображение
\[
\alpha:V\to K.
\]
Иными словами, для любых \(x,y\in V\) и \(\lambda\in K\)
\[
\alpha(x+y)=\alpha(x)+\alpha(y),
\qquad
\alpha(\lambda x)=\lambda\alpha(x).
\]
\end{definition}

Если в базисе \(e_1,\dots,e_n\) вектор имеет координаты
\[
x=x_1e_1+\cdots+x_ne_n,
\]
то линейность даёт
\[
\alpha(x)=a_1x_1+\cdots+a_nx_n,
\qquad
 a_i=\alpha(e_i).
\]
Поэтому функционал полностью определяется своими значениями на базисных векторах. В координатах он действует как строка коэффициентов, умноженная на столбец координат:
\[
\alpha(x)=
\begin{pmatrix}
 a_1&\cdots&a_n
\end{pmatrix}
\begin{pmatrix}
 x_1\\ \vdots\\ x_n
\end{pmatrix}.
\]

У Винберга среди основных примеров встречаются значение функции в точке, производная в точке, определённый интеграл и след матрицы. Все эти операции превращают сложный объект в одно число и линейно зависят от входа.

\begin{example}
На пространстве многочленов степени не выше двух функционал
\[
\alpha(f)=f(1)
\]
линеен. Если
\[
f(t)=a+bt+ct^2,
\]
то
\[
\alpha(f)=a+b+c.
\]
В базисе \(1,t,t^2\) функционалу соответствует строка
\[
\begin{pmatrix}1&1&1\end{pmatrix}.
\]
\end{example}

\subsection{Сопряжённое пространство и двойственный базис}

\begin{definition}
Множество всех линейных функционалов на \(V\) с обычными операциями сложения и умножения на скаляр называется сопряжённым, или двойственным, пространством и обозначается через \(V^*\).
\end{definition}

Пусть \(e_1,\dots,e_n\) --- базис пространства \(V\). Определим функционалы \(\varepsilon^1,\dots,\varepsilon^n\) условиями
\[
\varepsilon^i(e_j)=\delta_{ij},
\]
где \(\delta_{ij}=1\) при \(i=j\) и \(0\) иначе.

\begin{definition}
Система \(\varepsilon^1,\dots,\varepsilon^n\) называется базисом, двойственным базису \(e_1,\dots,e_n\).
\end{definition}

Функционал \(\varepsilon^i\) извлекает \(i\)-ю координату:
\[
\varepsilon^i(x)=x_i.
\]
Каждый функционал представляется единственным образом:
\[
\alpha=a_1\varepsilon^1+\cdots+a_n\varepsilon^n,
\qquad
 a_i=\alpha(e_i).
\]
Следовательно,
\[
\dim V^*=\dim V.
\]
Равенство размерностей не означает, что вектор и функционал являются одним объектом. Вектор задаёт направление или состояние, а функционал задаёт правило измерения этого состояния.

При смене базиса их координаты преобразуются по-разному. Если
\[
(e'_1,\dots,e'_n)=(e_1,\dots,e_n)C,
\]
то
\[
[x]_e=C[x]_{e'}.
\]
Если \(a_e\) --- строка коэффициентов функционала в старом базисе, то
\[
\alpha(x)=a_e[x]_e=a_eC[x]_{e'},
\]
поэтому
\[
a_{e'}=a_eC.
\]
Именно поэтому функционалы часто называют ковекторами.

\subsection{Второе сопряжённое пространство}

Каждый вектор \(x\in V\) можно рассматривать как функционал на \(V^*\):
\[
J(x)(\alpha)=\alpha(x).
\]

\begin{theorem}
Для конечномерного пространства отображение
\[
J:V\to V^{**},
\qquad
x\mapsto J(x),
\]
является изоморфизмом.
\end{theorem}

Этот изоморфизм не требует выбора базиса: вектор определяется тем, какие значения на нём принимают все линейные измерения. Для бесконечномерных пространств такое утверждение без дополнительных условий уже неверно.

\subsection{Ядро функционала и гиперплоскость}

Для ненулевого функционала \(\alpha\) множество
\[
\ker\alpha=\{x\in V:\alpha(x)=0\}
\]
является подпространством размерности \(\dim V-1\). Такое подпространство называется гиперплоскостью.

Геометрически уравнение
\[
\alpha(x)=0
\]
выбирает все направления, невидимые данному измерению. Неоднородное уравнение
\[
\alpha(x)=b
\]
задаёт параллельный сдвиг этой гиперплоскости.

\subsection{Аннулятор подпространства}

\begin{definition}
Аннулятором подпространства \(U\subseteq V\) называется подпространство
\[
U^\circ=
\{\alpha\in V^*:\alpha(u)=0\text{ для всех }u\in U\}.
\]
\end{definition}

Аннулятор состоит из всех линейных измерений, которые не различают движения внутри \(U\).

\begin{theorem}
Если \(V\) конечномерно, то
\[
\dim U^\circ=\dim V-\dim U.
\]
Кроме того,
\[
(U^\circ)^\circ=U.
\]
\end{theorem}

Если строки матрицы системы \(Ax=0\) понимать как функционалы \(\alpha_1,\dots,\alpha_m\), то пространство решений имеет вид
\[
\ker A
=
\{x:\alpha_i(x)=0\text{ для всех }i\}
=
\operatorname{span}(\alpha_1,\dots,\alpha_m)^\circ.
\]
Так язык сопряжённого пространства связывает систему уравнений с геометрией ограничений.

\paragraph{Связь с оптимизацией и машинным обучением.}

Для дифференцируемой скалярной функции \(F\) первая вариация в точке \(x\) действует на малое изменение \(h\) линейно:
\[
dF_x(h).
\]
Следовательно, \(dF_x\in V^*\). В координатах этот функционал записывается через градиент, но сам дифференциал по смыслу является ковектором: он сообщает, как изменение параметров влияет на целевое значение. Эта точка зрения понадобится при матричном дифференцировании, методе сопряжённых переменных и обратном распространении ошибки.

\subsection*{Задачи}

\begin{problem}
\textit{Задача Винберга.}
Пусть \(\dim V=n\). Показать, что функционалы
\[
\varepsilon^1,\dots,\varepsilon^n\in V^*
\]
образуют базис \(V^*\) тогда и только тогда, когда не существует ненулевого вектора \(x\in V\), для которого
\[
\varepsilon^1(x)=\cdots=\varepsilon^n(x)=0.
\]
\end{problem}

\begin{proof}[Решение]
Предположим сначала, что функционалы образуют базис. Если все они обращаются в нуль на \(x\), то любой функционал
\[
\alpha=c_1\varepsilon^1+\cdots+c_n\varepsilon^n
\]
также удовлетворяет \(\alpha(x)=0\). Для ненулевого \(x\) можно дополнить \(x\) до базиса пространства и построить функционал, равный единице на \(x\) и нулю на остальных базисных векторах. Получается противоречие, поэтому \(x=0\).

Обратно, рассмотрим отображение
\[
T:V\to K^n,
\qquad
T(x)=
\begin{pmatrix}
\varepsilon^1(x)\\ \vdots\\ \varepsilon^n(x)
\end{pmatrix}.
\]
Условие задачи означает, что \(\ker T=\{0\}\), то есть \(T\) инъективно. Пространства \(V\) и \(K^n\) имеют одинаковую размерность, поэтому \(T\) биективно. Для каждого \(j\) существует вектор \(x_j\), на котором
\[
\varepsilon^i(x_j)=\delta_{ij}.
\]
Если
\[
c_1\varepsilon^1+\cdots+c_n\varepsilon^n=0,
\]
то подстановка \(x_j\) даёт \(c_j=0\). Функционалы линейно независимы, а их число равно \(n=\dim V^*\), следовательно, они образуют базис.
\end{proof}

\begin{problem}
\textit{Задача Винберга.}
Пусть \(P_n\) --- пространство многочленов степени не выше \(n\), а \(x_0,\dots,x_n\) --- различные элементы поля. Для
\[
\varepsilon^i(f)=f(x_i)
\]
найти базис пространства \(P_n\), двойственный функционалам \(\varepsilon^0,\dots,\varepsilon^n\), и получить интерполяционную формулу Лагранжа.
\end{problem}

\begin{proof}[Решение]
Для каждого \(j\) определим многочлен
\[
L_j(t)=
\prod_{\substack{0\leq k\leq n\\k\neq j}}
\frac{t-x_k}{x_j-x_k}.
\]
При \(t=x_i\) получаем
\[
L_j(x_i)=\delta_{ij}.
\]
Следовательно, функционалы вычисления значения и многочлены \(L_j\) образуют двойственные базисы:
\[
\varepsilon^i(L_j)=\delta_{ij}.
\]
Разложение произвольного \(f\in P_n\) по базису \(L_0,\dots,L_n\) имеет вид
\[
f(t)=\sum_{j=0}^n \varepsilon^j(f)L_j(t)
=\sum_{j=0}^n f(x_j)
\prod_{\substack{0\leq k\leq n\\k\neq j}}
\frac{t-x_k}{x_j-x_k}.
\]
Это и есть интерполяционная формула Лагранжа. Она восстанавливает многочлен по \(n+1\) измерениям его значений.
\end{proof}

\begin{problem}
\textit{Практическая задача: восстановление состояния по датчикам.}
Состояние помещения описывается нормированными отклонениями
\[
x=(\tau,h,c)\in\R^3,
\]
где \(\tau\) --- температура, \(h\) --- влажность, \(c\) --- концентрация углекислого газа. Три датчика выдают
\[
\alpha_1(x)=\tau+h,
\qquad
\alpha_2(x)=h+c,
\qquad
\alpha_3(x)=\tau+c.
\]
Определить, позволяют ли показания однозначно восстановить состояние. Восстановить \(x\), если показания равны \(5,7,6\).
\end{problem}

\begin{proof}[Решение]
Коэффициенты функционалов образуют матрицу
\[
M=
\begin{pmatrix}
1&1&0\\
0&1&1\\
1&0&1
\end{pmatrix},
\qquad
\det M=2\neq0.
\]
Поэтому функционалы \(\alpha_1,\alpha_2,\alpha_3\) образуют базис двойственного пространства: никакое ненулевое изменение состояния не остаётся невидимым всем трём датчикам.

Из показаний получаем систему
\[
\tau+h=5,
\qquad
h+c=7,
\qquad
\tau+c=6.
\]
Складывая первое и третье уравнения и вычитая второе, находим
\[
2\tau=4,
\qquad
\tau=2.
\]
Затем
\[
h=3,
\qquad
c=4.
\]
Итак,
\[
x=(2,3,4).
\]
Линейная независимость датчиков здесь означает полную наблюдаемость трёхкомпонентного состояния.
\end{proof}

\begin{problem}
\textit{Практическая задача: нагрузка, не совершающая работу.}
Нормированный малый сдвиг механизма описывается вектором
\[
q=(u,v,\theta)\in\R^3,
\]
где \(u,v\) --- поступательные перемещения, а \(\theta\) --- малый поворот. Допустимые движения образуют подпространство
\[
U=\operatorname{span}
\left\{
(1,0,1),
(0,1,-1)
\right\}.
\]
Обобщённая нагрузка задаётся функционалом виртуальной работы
\[
\ell(q)=F_xu+F_yv+M\theta.
\]
Найти все нагрузки, которые не совершают работу ни на одном допустимом движении.
\end{problem}

\begin{proof}[Решение]
Требуется найти аннулятор \(U^\circ\). Функционал должен обращаться в нуль на обоих порождающих движениях:
\[
F_x+M=0,
\qquad
F_y-M=0.
\]
Полагая \(M=\lambda\), получаем
\[
(F_x,F_y,M)=\lambda(-1,1,1).
\]
Следовательно,
\[
U^\circ
=
\operatorname{span}
\left\{\ell_0\right\},
\qquad
\ell_0(u,v,\theta)=-u+v+\theta.
\]
Любая нагрузка, пропорциональная \(\ell_0\), ортогональна в смысле виртуальной работы всем допустимым движениям механизма. Размерность аннулятора равна
\[
3-\dim U=1,
\]
что согласуется с общим утверждением.
\end{proof}

\section{Билинейные и квадратичные формы}
\label{sec:bilinear-quadratic-forms}

Билинейная форма измеряет совместное действие двух векторов, причём по каждому аргументу зависимость остаётся линейной. Квадратичная форма получается, когда в оба аргумента подставляют один и тот же вектор. Эти конструкции описывают скалярные произведения, энергии, локальную кривизну функций потерь и устойчивость равновесий.

\subsection{Билинейная форма и её матрица}

\begin{definition}
Билинейной формой на векторном пространстве $V$ над полем $K$ называется отображение
\[
\alpha\colon V\times V\to K,
\]
линейное по каждому аргументу. Для любых $x_1,x_2,y_1,y_2\in V$ и $\lambda,\mu\in K$ выполняются, в частности,
\[
\alpha(\lambda x_1+\mu x_2,y)
=\lambda\alpha(x_1,y)+\mu\alpha(x_2,y),
\]
\[
\alpha(x,\lambda y_1+\mu y_2)
=\lambda\alpha(x,y_1)+\mu\alpha(x,y_2).
\]
\end{definition}

Пусть $e_1,\dots,e_n$ --- базис $V$, а $x$ и $y$ имеют координатные столбцы $X$ и $Y$. Числа
\[
a_{ij}=\alpha(e_i,e_j)
\]
образуют матрицу билинейной формы $A=(a_{ij})$, и
\[
\alpha(x,y)=X^{\trans}AY.
\]
Матрица полностью задаёт форму в выбранных координатах.

При замене базиса
\[
(e'_1,\dots,e'_n)=(e_1,\dots,e_n)C
\]
координаты преобразуются по правилу $X=CX'$ и $Y=CY'$. Поэтому новая матрица формы равна
\[
A'=C^{\trans}AC.
\]
Такое преобразование называется преобразованием по конгруэнтности. Его нельзя путать с формулой $C^{-1}AC$ для матрицы линейного оператора: у билинейной формы одновременно меняются оба аргумента.

\begin{definition}
Правым ядром билинейной формы называется подпространство
\[
\ker_{R}\alpha
=
\{y\in V\mid \alpha(x,y)=0\text{ для всех }x\in V\}.
\]
Левым ядром называется подпространство
\[
\ker_{L}\alpha
=
\{x\in V\mid \alpha(x,y)=0\text{ для всех }y\in V\}.
\]
Форма называется невырожденной, если оба ядра нулевые.
\end{definition}

В координатах условия имеют вид
\[
y\in\ker_R\alpha\iff AY=0,
\qquad
x\in\ker_L\alpha\iff A^{\trans}X=0.
\]
Поскольку квадратные матрицы $A$ и $A^{\trans}$ имеют одинаковый ранг,
\[
\dim\ker_R\alpha
=
\dim\ker_L\alpha
=
n-\operatorname{rank}A.
\]
Поэтому в конечномерном случае достаточно проверить нулевое значение одного из двух ядер. Для симметрической формы они совпадают, и обозначение $\ker\alpha$ становится однозначным. В частности, форма невырожденна тогда и только тогда, когда её матрица обратима.

\subsection{Симметрические и кососимметрические формы}

\begin{definition}
Билинейная форма называется симметрической, если
\[
\alpha(x,y)=\alpha(y,x),
\]
и кососимметрической, если
\[
\alpha(x,y)=-\alpha(y,x)
\]
для любых $x,y\in V$.
\end{definition}

В матричной записи эти условия имеют вид
\[
A^{\trans}=A
\qquad\text{и соответственно}\qquad
A^{\trans}=-A.
\]
Для кососимметрической формы над полем характеристики, отличной от двух,
\[
\alpha(x,x)=0
\]
для любого $x$. Такая форма измеряет не длину, а ориентированное попарное взаимодействие. В подходящем базисе её матрица состоит из блоков
\[
\begin{pmatrix}
0&1\\
-1&0
\end{pmatrix}
\]
и нулевого блока; отсюда, в частности, следует чётность её ранга.

\subsection{Квадратичная форма и поляризация}

\begin{definition}
Пусть $\alpha$ --- симметрическая билинейная форма над полем характеристики, отличной от двух. Функция
\[
q(x)=\alpha(x,x)
\]
называется квадратичной формой, ассоциированной с $\alpha$.
\end{definition}

Если $A$ --- матрица $\alpha$, то
\[
q(x)=X^{\trans}AX.
\]
Для двух переменных симметричная матрица
\[
A=
\begin{pmatrix}
a&b\\
b&c
\end{pmatrix}
\]
задаёт форму
\[
q(x_1,x_2)=ax_1^2+2bx_1x_2+cx_2^2.
\]
Коэффициент смешанного одночлена делится между двумя симметричными позициями матрицы. Например, форме
\[
q(x_1,x_2)=3x_1^2+4x_1x_2+5x_2^2
\]
соответствует матрица
\[
\begin{pmatrix}
3&2\\
2&5
\end{pmatrix}.
\]

Симметрическая билинейная форма восстанавливается по квадратичной форме формулой поляризации
\[
\alpha(x,y)
=
\frac12\bigl(q(x+y)-q(x)-q(y)\bigr).
\]
Следовательно, симметрические билинейные и квадратичные формы содержат одну и ту же информацию, записанную двумя способами.

\subsection{Диагональный и нормальный виды}

Смешанные члены затрудняют чтение формы. Подходящая замена базиса устраняет их.

\begin{theorem}
Для любой симметрической билинейной формы на конечномерном пространстве над полем характеристики, отличной от двух, существует базис, в котором её матрица диагональна.
\end{theorem}

В таком базисе
\[
q(x)=a_1x_1^2+\cdots+a_nx_n^2.
\]
Практически диагональный вид можно искать последовательным выделением квадратов. Например,
\[
2x^2+4xy+5y^2
=
2(x+y)^2+3y^2.
\]
Замена координат $u=x+y$, $v=y$ устраняет смешанный член и сразу показывает, что форма положительна на всех ненулевых векторах.

Для вещественной квадратичной формы ненулевые диагональные коэффициенты можно нормировать до $1$ и $-1$. Поэтому в некотором базисе
\[
q(x)
=
x_1^2+\cdots+x_k^2
-x_{k+1}^2-\cdots-x_{k+l}^2.
\]
Оставшиеся координаты, если они есть, входят с нулевыми коэффициентами.

\begin{theorem}[закон инерции]
Числа $k$ и $l$ в нормальном виде вещественной квадратичной формы не зависят от выбора базиса.
\end{theorem}

Пара $(k,l)$ называется сигнатурой формы. Число нулевых направлений равно
\[
n-k-l=\dim\ker q.
\]
Положительные направления увеличивают значение формы, отрицательные уменьшают его, а нулевые не обнаруживаются квадратичным членом.

\subsection{Положительная определённость}

\begin{definition}
Вещественная квадратичная форма $q$ называется положительно определённой, если
\[
q(x)>0
\]
для каждого ненулевого $x$. Она называется отрицательно определённой, если $q(x)<0$ при $x\neq0$, и неопределённой, если принимает значения обоих знаков.
\end{definition}

Положительно определённая форма имеет сигнатуру $(n,0)$. В механике она задаёт энергию устойчивого малого отклонения; в оптимизации --- положительную локальную кривизну во всех направлениях.

Для симметрической матрицы $A\in\R^{n\times n}$ обозначим через
\[
\Delta_k
=
\det A_{1:k,1:k}
\]
определитель её левого верхнего блока порядка $k$.

\begin{criterion}[критерий Сильвестра]
Квадратичная форма $q(x)=X^{\trans}AX$ положительно определена тогда и только тогда, когда
\[
\Delta_1>0,\quad \Delta_2>0,\quad\dots,\quad\Delta_n>0.
\]
\end{criterion}

Для матрицы второго порядка
\[
A=
\begin{pmatrix}
a&b\\b&c\end{pmatrix}
\]
критерий принимает вид
\[
a>0,
\qquad
ac-b^2>0.
\]
Одной положительности диагональных элементов недостаточно: смешанное взаимодействие может сделать форму неопределённой.

\begin{remark}
Критерий Сильвестра в приведённой форме проверяет строгую положительную определённость. Для положительной полуопределённости проверки только угловых миноров недостаточно; требуется анализ всех главных миноров или другой эквивалентный критерий.
\end{remark}

\subsection{Квадратичная форма как локальная модель}

Пусть гладкая функция $F$ зависит от вектора параметров $z$. В окрестности точки $z_0$ её приращение имеет вид
\[
F(z_0+h)
=
F(z_0)
+\nabla F(z_0)^{\trans}h
+\frac12 h^{\trans}H(z_0)h
+o(\lVert h\rVert^2),
\]
где $H(z_0)$ --- матрица вторых производных, или гессиан. В стационарной точке линейный член исчезает, и характер точки определяется квадратичной формой гессиана. Положительная определённость указывает на строгий локальный минимум, отрицательная --- на максимум, а неопределённость --- на седловую точку.

Эта связь будет использоваться в многомерном анализе, оптимизации и обучении моделей. Классификация локальной кривизны сводится к исследованию симметрической матрицы по конгруэнтности.

\subsection*{Задачи}

\begin{problem}
\textit{Задача Винберга.}
Доказать невырожденность билинейных форм
\[
\alpha(X,Y)=\operatorname{tr}(XY)
\]
на пространстве квадратных матриц $L_n(K)$ и
\[
\beta(x,y)=x_1y_2-x_2y_1
\]
на пространстве $K^2$.
\end{problem}

\begin{proof}[Решение]
Пусть $Y\neq0$. Тогда существует ненулевой элемент $y_{ij}$. Возьмём матричную единицу $X=E_{ji}$. Получаем
\[
\alpha(E_{ji},Y)=\operatorname{tr}(E_{ji}Y)=y_{ij}\neq0.
\]
Следовательно, ненулевой $Y$ не принадлежит ядру, и форма $\alpha$ невырожденна.

Теперь пусть $y=(y_1,y_2)\neq0$. Если $y_1\neq0$, выберем $x=(0,-1)$; тогда
\[
\beta(x,y)=y_1\neq0.
\]
Если $y_1=0$, то $y_2\neq0$, и при $x=(1,0)$ имеем
\[
\beta(x,y)=y_2\neq0.
\]
Значит, форма $\beta$ также невырожденна. Геометрически она вычисляет ориентированную площадь параллелограмма, построенного на двух векторах.
\end{proof}

\begin{problem}
\textit{Задача Винберга.}
Найти сигнатуру вещественной симметрической билинейной формы
\[
\alpha(X,Y)=\operatorname{tr}(XY)
\]
на пространстве $L_n(\R)$.
\end{problem}

\begin{proof}[Решение]
Любая вещественная матрица единственным образом представляется как
\[
X=S+K,
\qquad
S^{\trans}=S,
\qquad
K^{\trans}=-K.
\]
Подпространства симметрических и кососимметрических матриц ортогональны относительно $\alpha$, поскольку
\[
\operatorname{tr}(SK)
=
\operatorname{tr}((SK)^{\trans})
=
-\operatorname{tr}(KS)
=
-\operatorname{tr}(SK).
\]
Следовательно, $\operatorname{tr}(SK)=0$.

Для симметрической матрицы
\[
q(S)=\operatorname{tr}(S^2)
=
\sum_i s_{ii}^2+2\sum_{i<j}s_{ij}^2>0
\]
при $S\neq0$. На пространстве кососимметрических матриц
\[
q(K)=\operatorname{tr}(K^2)
=
-2\sum_{i<j}k_{ij}^2<0
\]
при $K\neq0$.

Размерности этих подпространств равны
\[
\dim\operatorname{Sym}_n
=\frac{n(n+1)}2,
\qquad
\dim\operatorname{Skew}_n
=\frac{n(n-1)}2.
\]
Поэтому сигнатура формы равна
\[
\left(
\frac{n(n+1)}2,
\frac{n(n-1)}2
\right).
\]
\end{proof}

\begin{problem}
\textit{Практическая задача: устойчивость двух узлов панели.}
Два узла гибкой панели могут совершать малые вертикальные перемещения $u$ и $v$. Каждый узел соединён с неподвижной опорой пружиной единичной жёсткости, а узлы соединены между собой ещё одной такой пружиной. Потенциальная энергия равна
\[
E(u,v)
=
\frac12u^2+\frac12v^2+\frac12(u-v)^2.
\]
Проверить устойчивость положения $u=v=0$.
\end{problem}

\begin{proof}[Решение]
После упрощения
\[
E(u,v)=u^2-uv+v^2
=
\begin{pmatrix}u&v\end{pmatrix}
\begin{pmatrix}
1&-\tfrac12\\
-\tfrac12&1
\end{pmatrix}
\begin{pmatrix}u\\v\end{pmatrix}.
\]
Угловые миноры матрицы жёсткости равны
\[
\Delta_1=1,
\qquad
\Delta_2=1-\frac14=\frac34.
\]
Оба положительны, поэтому энергия положительно определена. Эквивалентно,
\[
E(u,v)
=
\left(u-\frac v2\right)^2+\frac34v^2>0
\]
при $(u,v)\neq(0,0)$. Любое малое ненулевое отклонение увеличивает энергию, следовательно, положение равновесия устойчиво.
\end{proof}

\begin{problem}
\textit{Практическая задача: седловое направление в параметрах конструкции.}
Пусть $a$ --- нормированное изменение толщины элемента, а $b$ --- нормированное изменение его кривизны. Вблизи текущей конструкции изменение целевой функции приближённо равно
\[
\Delta J(a,b)=2a^2+4ab+b^2.
\]
Меньшее значение $J$ считается лучшим. Определить характер текущей конструкции по этой квадратичной модели.
\end{problem}

\begin{proof}[Решение]
Матрица формы равна
\[
H=
\begin{pmatrix}
2&2\\
2&1
\end{pmatrix}.
\]
Хотя оба диагональных элемента положительны,
\[
\det H=2-4=-2<0.
\]
Следовательно, форма неопределённа. Действительно, вдоль координатных направлений
\[
\Delta J(a,0)=2a^2>0,
\qquad
\Delta J(0,b)=b^2>0,
\]
но для согласованного изменения $b=-2a$
\[
\Delta J(a,-2a)
=2a^2-8a^2+4a^2
=-2a^2<0.
\]
По отдельности изменение толщины или кривизны ухудшает цель, однако их совместное изменение в правильной пропорции улучшает её. Текущая конструкция является седловой точкой локальной квадратичной модели, а не локальным минимумом.
\end{proof}

\section{Евклидовы и эрмитовы пространства}
\label{sec:euclidean-hermitian}

Квадратичная форма позволяет измерять величину вектора. Если эта форма положительно определена, она задаёт скалярное произведение, а вместе с ним длины, углы, ортогональность и расстояния. Именно здесь линейная алгебра превращается в геометрию и даёт строгую основу проекциям, наименьшим квадратам и ортогональным разложениям.

\subsection{Скалярное произведение и норма}

\begin{definition}
Евклидовым пространством называется вещественное векторное пространство \(V\), на котором задано положительно определённое симметрическое билинейное отображение
\[
(\,\cdot\,,\,\cdot\,)\colon V\times V\to\R.
\]
Оно называется скалярным произведением.
\end{definition}

Положительная определённость означает
\[
(x,x)>0
\quad\text{при }x\neq0.
\]
Поэтому величина
\[
\|x\|=\sqrt{(x,x)}
\]
корректно определяет длину вектора. В стандартном пространстве \(\R^n\)
\[
(x,y)=x^{\trans}y
=
\sum_{i=1}^n x_i y_i.
\]

\begin{theorem}[неравенство Коши--Буняковского]
Для любых \(x,y\in V\)
\[
|(x,y)|\leq \|x\|\,\|y\|.
\]
Равенство достигается тогда и только тогда, когда векторы линейно зависимы.
\end{theorem}

Это неравенство позволяет определить угол между ненулевыми векторами:
\[
\cos\theta=\frac{(x,y)}{\|x\|\,\|y\|}.
\]
Расстояние задаётся формулой
\[
\rho(x,y)=\|x-y\|.
\]

\subsection{Ортогональность и ортогональное дополнение}

\begin{definition}
Векторы \(x\) и \(y\) называются ортогональными, если
\[
(x,y)=0.
\]
Для подпространства \(U\subseteq V\) его ортогональным дополнением называется
\[
U^\perp
=
\{z\in V:(z,u)=0\text{ для всех }u\in U\}.
\]
\end{definition}

\begin{theorem}
Для любого подпространства \(U\) конечномерного евклидова пространства
\[
V=U\oplus U^\perp.
\]
Следовательно, каждый \(x\in V\) единственным образом представляется как
\[
x=y+z,
\qquad
y\in U,
\quad
z\in U^\perp.
\]
\end{theorem}

Вектор \(y\) называется ортогональной проекцией \(x\) на \(U\), а \(z=x-y\) --- ортогональной составляющей. Теорема утверждает не только существование разложения, но и его единственность.

Если \(e_1,\dots,e_k\) --- ортонормированный базис \(U\), то
\[
\operatorname{pr}_U x
=
\sum_{i=1}^k (x,e_i)e_i.
\]
Для ортогонального, но не нормированного базиса
\[
\operatorname{pr}_U x
=
\sum_{i=1}^k
\frac{(x,e_i)}{(e_i,e_i)}e_i.
\]

\begin{theorem}[о ближайшей точке]
Ортогональная проекция \(\operatorname{pr}_U x\) является единственным вектором подпространства \(U\), ближайшим к \(x\):
\[
\|x-\operatorname{pr}_U x\|
=
\min_{u\in U}\|x-u\|.
\]
\end{theorem}

Практический смысл прост: если точное представление объекта в допустимом подпространстве невозможно, ортогональная проекция даёт наилучшее приближение в евклидовой метрике.

\subsection{Ортонормированные базисы и процесс Грама--Шмидта}

\begin{definition}
Система \(q_1,\dots,q_k\) называется ортонормированной, если
\[
(q_i,q_j)=\delta_{ij}.
\]
\end{definition}

Пусть \(a_1,\dots,a_k\) линейно независимы. Процесс Грама--Шмидта строит из них ортогональную систему с той же линейной оболочкой:
\[
u_1=a_1,
\]
\[
u_j
=
a_j-
\sum_{i=1}^{j-1}
\frac{(a_j,u_i)}{(u_i,u_i)}u_i,
\qquad j\geq2.
\]
После нормировки
\[
q_j=\frac{u_j}{\|u_j\|}
\]
получается ортонормированная система.

На каждом шаге из нового вектора вычитаются его проекции на уже построенные направления. Остаток ортогонален всем предыдущим векторам и содержит только действительно новое направление.

\begin{computationalnote}
В вычислениях классическая запись Грама--Шмидта может накапливать ошибку округления. Для численной реализации используют модифицированный процесс Грама--Шмидта или отражения Хаусхолдера. Теоретическая идея остаётся той же: отделить взаимно ортогональные направления.
\end{computationalnote}

\subsection{Ортогональные матрицы и QR-разложение}

\begin{definition}
Квадратная вещественная матрица \(Q\) называется ортогональной, если
\[
Q^{\trans}Q=I.
\]
\end{definition}

Столбцы ортогональной матрицы образуют ортонормированный базис. Поэтому
\[
Q^{-1}=Q^{\trans},
\qquad
\|Qx\|=\|x\|.
\]
Ортогональные преобразования сохраняют длины, углы и расстояния.

Если к линейно независимым столбцам матрицы \(A\in\R^{m\times n}\) применить ортогонализацию, получим
\[
A=QR,
\]
где столбцы \(Q\in\R^{m\times n}\) ортонормированы, а \(R\in\R^{n\times n}\) верхнетреугольна. Это QR-разложение. Оно используется для устойчивого решения систем и задач наименьших квадратов.

\subsection{Наименьшие квадраты}

Пусть система
\[
Ax=b,
\qquad
A\in\R^{m\times n},
\]
несовместна. Тогда вектор \(b\) не лежит в пространстве столбцов \(A\). Вместо точного решения ищут параметры, для которых остаток
\[
r=b-Ax
\]
имеет минимальную норму.

\begin{definition}
Решением задачи наименьших квадратов называется вектор \(x_\ast\), минимизирующий
\[
\|Ax-b\|^2.
\]
\end{definition}

Ближайший к \(b\) достижимый вектор \(Ax_\ast\) является ортогональной проекцией \(b\) на \(\operatorname{im}A\). Поэтому остаток ортогонален каждому столбцу \(A\):
\[
A^{\trans}(b-Ax_\ast)=0.
\]
Отсюда следуют нормальные уравнения
\[
A^{\trans}A x_\ast=A^{\trans}b.
\]

Если столбцы \(A\) линейно независимы, то \(A^{\trans}A\) положительно определена и
\[
x_\ast
=
(A^{\trans}A)^{-1}A^{\trans}b.
\]
Матрица ортогональной проекции на пространство столбцов имеет вид
\[
P_A
=
A(A^{\trans}A)^{-1}A^{\trans},
\qquad
Ax_\ast=P_A b.
\]

\begin{computationalnote}
Нормальные уравнения удобны для вывода, но при плохой обусловленности усиливают численную ошибку. На практике задачу чаще решают через QR-разложение:
\[
A=QR,
\qquad
Rx_\ast=Q^{\trans}b.
\]
\end{computationalnote}

\subsection{Эрмитовы пространства}

Для комплексных векторов обычная билинейная форма \(x^{\trans}y\) не задаёт положительную длину. Комплексное сопряжение приводит к эрмитову скалярному произведению.

\begin{definition}
Эрмитовым пространством называется комплексное векторное пространство \(V\) с отображением
\[
\langle\,\cdot\,,\,\cdot\,\rangle\colon V\times V\to\C,
\]
которое сопряжённо-линейно по первому аргументу, линейно по второму, удовлетворяет
\[
\langle y,x\rangle
=
\overline{\langle x,y\rangle}
\]
и положительно определено:
\[
\langle x,x\rangle>0
\quad\text{при }x\neq0.
\]
\end{definition}

В \(\C^n\)
\[
\langle x,y\rangle
=
x^*y
=
\sum_{i=1}^n \overline{x_i}y_i,
\]
где
\[
A^*=\overline{A}^{\trans}
\]
--- сопряжённо-транспонированная матрица.

Все основные конструкции евклидова пространства сохраняются. Для ортонормированного базиса \(e_1,\dots,e_k\)
\[
\operatorname{pr}_U x
=
\sum_{i=1}^k
\langle e_i,x\rangle e_i.
\]
Комплексная матрица \(U\) называется унитарной, если
\[
U^*U=I.
\]
Унитарные преобразования сохраняют эрмитову норму и комплексные углы. В формулах наименьших квадратов транспонирование заменяется на сопряжённое транспонирование:
\[
A^*A x_\ast=A^*b.
\]

\subsection*{Задачи}

\begin{problem}
\textit{Задача Винберга.}
Применяя процесс ортогонализации к столбцам, доказать, что каждая матрица
\[
A\in GL_n(\R)
\]
единственным образом представляется в виде
\[
A=QR,
\]
где \(Q\) ортогональна, а \(R\) верхнетреугольна и имеет положительные диагональные элементы.
\end{problem}

\begin{proof}[Решение]
Обозначим столбцы \(A\) через \(a_1,\dots,a_n\). Они линейно независимы. После процесса Грама--Шмидта получаем ортонормированный базис \(q_1,\dots,q_n\), причём для каждого \(j\)
\[
a_j\in\operatorname{span}(q_1,\dots,q_j).
\]
Поэтому
\[
a_j=\sum_{i=1}^j r_{ij}q_i.
\]
Коэффициенты \(r_{ij}\) образуют верхнетреугольную матрицу \(R\), а матрица \(Q\) состоит из столбцов \(q_i\). Значит,
\[
A=QR.
\]
Выбирая знак \(q_j\) так, чтобы \(r_{jj}>0\), получаем требуемую диагональ.

Пусть теперь
\[
A=Q_1R_1=Q_2R_2
\]
--- два таких разложения. Тогда
\[
Q_2^{\trans}Q_1=R_2R_1^{-1}.
\]
Левая часть ортогональна, правая верхнетреугольна и имеет положительную диагональ. Единственная матрица с этими свойствами --- единичная. Следовательно,
\[
Q_1=Q_2,
\qquad
R_1=R_2.
\]
\end{proof}

\begin{problem}
\textit{Задача по материалу Винберга.}
В пространстве многочленов степени не выше трёх зададим скалярное произведение
\[
(f,g)=\int_0^1 f(x)g(x)\,dx.
\]
Найти квадратный трёхчлен, ближайший к \(x^3\), и расстояние до него.
\end{problem}

\begin{proof}[Решение]
Подпространство квадратных трёхчленов равно
\[
U=\operatorname{span}(1,x,x^2).
\]
После ортогонализации системы \(1,x,x^2,x^3\) последний ортогональный вектор имеет вид
\[
u_4
=
x^3-\frac32x^2+\frac35x-\frac1{20}.
\]
Он ортогонален всему \(U\). Следовательно,
\[
x^3
=
\left(
\frac32x^2-\frac35x+\frac1{20}
\right)
+
u_4,
\]
и ближайшим квадратным трёхчленом является
\[
p_\ast(x)
=
\frac32x^2-\frac35x+\frac1{20}.
\]
Квадрат расстояния равен
\[
\|u_4\|^2
=
\int_0^1 u_4(x)^2\,dx
=
\frac1{2800}.
\]
Поэтому
\[
\|x^3-p_\ast\|
=
\frac1{20\sqrt7}.
\]
\end{proof}

\begin{problem}
\textit{Практическая задача: калибровка датчика температуры.}
Пусть \(u\) --- нормированное выходное напряжение датчика, а \(T\) --- температура по эталонному термометру. Получены измерения
\[
(-1,-1),\qquad (0,1),\qquad (1,2).
\]
Найти линейную калибровку
\[
T\approx au+b
\]
методом наименьших квадратов.
\end{problem}

\begin{proof}[Решение]
Матрица модели и вектор измерений равны
\[
A=
\begin{pmatrix}
-1&1\\
0&1\\
1&1
\end{pmatrix},
\qquad
t=
\begin{pmatrix}
-1\\1\\2
\end{pmatrix}.
\]
Нормальные уравнения имеют вид
\[
A^{\trans}A
\begin{pmatrix}a\\b\end{pmatrix}
=
A^{\trans}t.
\]
Вычисляя,
\[
A^{\trans}A
=
\begin{pmatrix}
2&0\\
0&3
\end{pmatrix},
\qquad
A^{\trans}t
=
\begin{pmatrix}
3\\2
\end{pmatrix}.
\]
Следовательно,
\[
a=\frac32,
\qquad
b=\frac23.
\]
Искомая калибровка:
\[
T\approx\frac32u+\frac23.
\]
Предсказания равны
\[
-\frac56,\qquad \frac23,\qquad \frac{13}{6}.
\]
Остаток
\[
r=t-A
\begin{pmatrix}a\\b\end{pmatrix}
=
\begin{pmatrix}
-\frac16\\[2mm]
\frac13\\[2mm]
-\frac16
\end{pmatrix}
\]
не равен нулю, потому что три шумных измерения не лежат точно на одной прямой. При этом
\[
A^{\trans}r=0,
\]
то есть ошибка ортогональна обоим направлениям модели.
\end{proof}

\begin{problem}
\textit{Практическая задача: приближение формы двумя приводами.}
Два привода создают в трёх контрольных точках панели профили перемещений
\[
q_1=
\begin{pmatrix}
1\\0\\1
\end{pmatrix},
\qquad
q_2=
\begin{pmatrix}
0\\1\\1
\end{pmatrix}.
\]
Требуемый профиль равен
\[
d=
\begin{pmatrix}
1\\1\\3
\end{pmatrix}.
\]
Найти команды приводов, дающие ближайший достижимый профиль.
\end{problem}

\begin{proof}[Решение]
Запишем
\[
A=
\begin{pmatrix}
1&0\\
0&1\\
1&1
\end{pmatrix},
\qquad
c=
\begin{pmatrix}
c_1\\c_2
\end{pmatrix}.
\]
Точный профиль \(Ac=d\) недостижим: первые две строки требуют
\(c_1=c_2=1\), тогда третья компонента равна \(2\), а не \(3\).

Нормальные уравнения:
\[
A^{\trans}A c=A^{\trans}d,
\]
\[
\begin{pmatrix}
2&1\\
1&2
\end{pmatrix}
c
=
\begin{pmatrix}
4\\4
\end{pmatrix}.
\]
Отсюда
\[
c_1=c_2=\frac43.
\]
Ближайший профиль:
\[
Ac_\ast
=
\begin{pmatrix}
\frac43\\[1mm]
\frac43\\[1mm]
\frac83
\end{pmatrix}.
\]
Остаток
\[
d-Ac_\ast
=
\begin{pmatrix}
-\frac13\\[1mm]
-\frac13\\[1mm]
\frac13
\end{pmatrix}
\]
ортогонален обоим профилям приводов:
\[
q_1^{\trans}(d-Ac_\ast)=0,
\qquad
q_2^{\trans}(d-Ac_\ast)=0.
\]
Иными словами, оставшуюся ошибку уже невозможно уменьшить никаким изменением доступных команд.
\end{proof}

\section{Собственные значения, диагонализация и спектральная теорема}
\label{sec:eigenvalues-spectral-theorem}

Матрица линейного оператора зависит от базиса, но сам оператор от выбора координат не зависит. Поэтому естественно искать направления, в которых действие оператора не смешивает координаты, а лишь умножает вектор на число. Такие направления позволяют разложить связанную систему на независимые режимы: деформации --- на собственные формы, теплоперенос --- на затухающие моды, а повторное применение преобразования --- на отдельные геометрические прогрессии.

\subsection{Собственные значения и собственные подпространства}

Пусть $A\colon V\to V$ --- линейный оператор в конечномерном пространстве над полем $K$.

\begin{definition}
Ненулевой вектор $v\in V$ называется собственным вектором оператора $A$, если существует число $\lambda\in K$, для которого
\[
Av=\lambda v.
\]
Число $\lambda$ называется соответствующим собственным значением.
\end{definition}

Геометрически прямая $\operatorname{span}(v)$ остаётся на месте: оператор может изменить длину вектора и, при отрицательном $\lambda$, его направление, но не выводит вектор из этой прямой. Нулевой вектор не считается собственным, поскольку равенство $A0=\lambda0$ выполняется при любом $\lambda$ и ничего не сообщает об операторе.

Для фиксированного $\lambda$ равенство $Av=\lambda v$ эквивалентно системе
\[
(A-\lambda I)v=0.
\]

\begin{definition}
Собственным подпространством, отвечающим значению $\lambda$, называется
\[
V_\lambda(A)=\ker(A-\lambda I).
\]
Оно состоит из нулевого вектора и всех собственных векторов со значением $\lambda$.
\end{definition}

Ненулевой собственный вектор существует тогда и только тогда, когда оператор $A-\lambda I$ вырожден. В координатах это условие имеет вид
\[
\det(A-\lambda I)=0.
\]

\begin{definition}
Характеристическим многочленом матрицы $A\in K^{n\times n}$ называется
\[
p_A(t)=\det(tI-A).
\]
Его корни являются собственными значениями матрицы.
\end{definition}

При смене базиса матрица оператора преобразуется по правилу $A'=C^{-1}AC$. Тогда
\[
\det(tI-A')
=
\det\bigl(C^{-1}(tI-A)C\bigr)
=
\det(tI-A),
\]
поэтому характеристический многочлен, собственные значения, след и определитель не зависят от базиса. При разложении
\[
p_A(t)=t^n-(\operatorname{tr}A)t^{n-1}+\cdots+(-1)^n\det A
\]
получаем, с учётом кратностей,
\[
\sum_{i=1}^{n}\lambda_i=\operatorname{tr}A,
\qquad
\prod_{i=1}^{n}\lambda_i=\det A.
\]

\begin{example}
Для матрицы
\[
A=
\begin{pmatrix}
2&1\\
1&2
\end{pmatrix}
\]
характеристический многочлен равен
\[
p_A(t)=
\det
\begin{pmatrix}
t-2&-1\\
-1&t-2
\end{pmatrix}
=(t-2)^2-1=(t-1)(t-3).
\]
При $\lambda=3$ система $(A-3I)v=0$ даёт направление $(1,1)$, а при $\lambda=1$ --- направление $(1,-1)$. Оператор усиливает совместную компоненту двух координат в три раза, а противоположную --- в один раз.
\end{example}

Над полем $\C$ характеристический многочлен всегда разлагается на линейные множители. Над $\R$ это может быть неверно. Например, поворот плоскости на угол, отличный от $0$ и $\pi$, не имеет вещественных собственных направлений; его собственные значения появляются после перехода к комплексному пространству.

\subsection{Независимость собственных направлений и кратности}

\begin{theorem}
Собственные подпространства, отвечающие различным собственным значениям, линейно независимы. В частности, собственные векторы, отвечающие попарно различным собственным значениям, линейно независимы.
\end{theorem}

\begin{proofidea}
Пусть $v_1+\cdots+v_k=0$, где $Av_i=\lambda_i v_i$ и значения $\lambda_i$ различны. Применение $A$ и вычитание исходного равенства, умноженного на $\lambda_k$, устраняет $v_k$ и оставляет зависимость между $k-1$ векторами. Индукция показывает, что все $v_i$ равны нулю.
\end{proofidea}

\begin{definition}
Алгебраической кратностью собственного значения $\lambda$ называется его кратность как корня характеристического многочлена. Геометрической кратностью называется
\[
g_\lambda=\dim V_\lambda(A)=\dim\ker(A-\lambda I).
\]
\end{definition}

Для каждого собственного значения
\[
1\leq g_\lambda\leq a_\lambda,
\]
где $a_\lambda$ --- алгебраическая кратность. Первая граница следует из существования собственного вектора, а вторая отражает то, что число независимых собственных направлений не может превышать кратность корня. Случай строгого неравенства $g_\lambda<a_\lambda$ является причиной невозможности обычной диагонализации и будет разобран при изучении жордановой формы.

\subsection{Диагонализация}

\begin{definition}
Оператор называется диагонализируемым, если в пространстве существует базис из его собственных векторов. Эквивалентно, матрица $A$ диагонализируема, если существует обратимая матрица $S$, такая что
\[
S^{-1}AS=\Lambda,
\]
где $\Lambda$ диагональна.
\end{definition}

Столбцы $S$ являются собственными векторами, а диагональные элементы $\Lambda$ --- соответствующими собственными значениями:
\[
AS=S\Lambda.
\]

\begin{theorem}
Пусть характеристический многочлен оператора $A$ разлагается над $K$ на линейные множители. Оператор $A$ диагонализируем тогда и только тогда, когда геометрическая кратность каждого собственного значения равна его алгебраической кратности:
\[
g_\lambda=a_\lambda.
\]
Эквивалентно,
\[
\sum_\lambda \dim V_\lambda(A)=n.
\]
\end{theorem}

Из теоремы сразу следует достаточное условие: матрица порядка $n$ с $n$ различными собственными значениями диагонализируема. Обратное неверно: единичная матрица диагональна, хотя имеет единственное собственное значение.

Диагонализация упрощает повторное применение оператора. Если
\[
A=S\Lambda S^{-1},
\]
то
\[
A^k=S\Lambda^kS^{-1},
\qquad
\Lambda^k=\operatorname{diag}(\lambda_1^k,\dots,\lambda_n^k).
\]
Более общо, для многочлена $f$
\[
f(A)=Sf(\Lambda)S^{-1}.
\]
Поэтому динамика итерации $x_{k+1}=Ax_k$ определяется модулями собственных значений. Компонента вдоль собственного вектора умножается после $k$ шагов на $\lambda^k$: при $|\lambda|<1$ она затухает, при $|\lambda|>1$ растёт, а при отрицательном $\lambda$ дополнительно меняет знак.

\subsection{Самосопряжённые операторы и спектральная теорема}

Пусть $V$ --- вещественное евклидово пространство со скалярным произведением $(\cdot,\cdot)$.

\begin{definition}
Оператор $A\colon V\to V$ называется самосопряжённым, если
\[
(Ax,y)=(x,Ay)
\]
для любых $x,y\in V$. В ортонормированном базисе его матрица симметрична:
\[
A^{\trans}=A.
\]
\end{definition}

Для произвольной матрицы наличие полного набора собственных векторов приходится проверять. Для симметрической матрицы это гарантировано структурой скалярного произведения.

\begin{theorem}[спектральная теорема]
Для любого самосопряжённого оператора в конечномерном вещественном евклидовом пространстве существует ортонормированный базис из собственных векторов. Все собственные значения вещественны, а собственные подпространства, отвечающие различным значениям, ортогональны.
\end{theorem}

Если $q_1,\dots,q_n$ --- ортонормированные собственные векторы и $Q=(q_1\ \cdots\ q_n)$, то
\[
Q^{\trans}Q=I,
\qquad
A=Q\Lambda Q^{\trans}.
\]
Эквивалентная форма спектрального разложения:
\[
A=\sum_{i=1}^{n}\lambda_i q_iq_i^{\trans}.
\]
Матрица $q_iq_i^{\trans}$ является ортогональным проектором на прямую $\operatorname{span}(q_i)$. Следовательно, спектральное разложение представляет оператор как сумму независимых действий по взаимно перпендикулярным направлениям.

Если
\[
x=\sum_{i=1}^{n}\alpha_iq_i,
\]
то
\[
Ax=\sum_{i=1}^{n}\lambda_i\alpha_iq_i,
\qquad
x^{\trans}Ax=\sum_{i=1}^{n}\lambda_i\alpha_i^2.
\]
Отсюда симметрическая матрица положительно определена тогда и только тогда, когда все её собственные значения положительны.

\begin{definition}
Для симметрической матрицы $A$ и ненулевого $x$ отношение
\[
\mathcal{R}_A(x)=\frac{x^{\trans}Ax}{x^{\trans}x}
\]
называется отношением Рэлея.
\end{definition}

В собственном базисе отношение Рэлея является взвешенным средним собственных значений:
\[
\mathcal{R}_A(x)
=
\frac{\sum_i\lambda_i\alpha_i^2}{\sum_i\alpha_i^2}.
\]
Поэтому
\[
\lambda_{\min}\leq \mathcal{R}_A(x)\leq \lambda_{\max}.
\]
Минимум и максимум достигаются на собственных векторах, отвечающих крайним собственным значениям. В механике эта величина связывает форму деформации с её энергией; в численных методах она позволяет оценивать крайние части спектра.

В комплексном эрмитовом пространстве утверждение сохраняется с заменой транспонирования на сопряжённое транспонирование: эрмитова матрица допускает разложение $A=Q\Lambda Q^*$ с унитарной $Q$ и вещественной $\Lambda$.

\subsection*{Задачи}

\begin{problem}
\textit{Задача Винберга.}
Доказать, что линейный оператор $P$ является проектором для некоторого разложения
\[
V=U\oplus W
\]
тогда и только тогда, когда
\[
P^2=P.
\]
\end{problem}

\begin{proof}[Решение]
Если $P$ --- проектор на $U$ параллельно $W$, то любой $x\in V$ единственным образом представляется как $x=u+w$, где $u\in U$, $w\in W$, и
\[
P(u+w)=u.
\]
Повторное применение ничего не меняет:
\[
P^2x=P(u)=u=Px.
\]

Обратно, пусть $P^2=P$. Положим
\[
U=\operatorname{im}P,
\qquad
W=\ker P.
\]
Для любого $x$ имеем разложение
\[
x=Px+(x-Px).
\]
Первое слагаемое принадлежит $U$, а второе --- $W$, поскольку
\[
P(x-Px)=Px-P^2x=0.
\]
Если $z\in U\cap W$, то $z=Py$ для некоторого $y$ и одновременно $Pz=0$. Но
\[
z=Py=P^2y=Pz=0.
\]
Следовательно, $V=U\oplus W$, причём $P$ действует как единица на $U$ и как нуль на $W$. Его спектр содержится в $\{0,1\}$, а базис, составленный из базисов $U$ и $W$, диагонализирует $P$.
\end{proof}

\begin{problem}
\textit{Задача Винберга.}
Над полем характеристики, отличной от $2$, доказать, что оператор $R$ является отражением относительно некоторого подпространства $U$ параллельно $W$ тогда и только тогда, когда
\[
R^2=I.
\]
\end{problem}

\begin{proof}[Решение]
Для отражения выполняется
\[
R(u+w)=u-w,
\qquad u\in U,
\quad w\in W,
\]
поэтому $R^2(u+w)=u+w$.

Пусть теперь $R^2=I$. Для любого $x\in V$ определим
\[
u=\frac{x+Rx}{2},
\qquad
w=\frac{x-Rx}{2}.
\]
Тогда $x=u+w$ и
\[
Ru=\frac{Rx+R^2x}{2}=u,
\qquad
Rw=\frac{Rx-R^2x}{2}=-w.
\]
Значит,
\[
U=\ker(R-I),
\qquad
W=\ker(R+I),
\]
и $V=U\oplus W$. В базисе из этих двух собственных подпространств матрица $R$ диагональна и содержит только значения $1$ и $-1$.
\end{proof}

\begin{problem}
\textit{Практическая задача: выравнивание температуры в двух зонах.}
Отклонения температур двух соседних зон от нормативного уровня после каждого шага обмена теплом обновляются по закону
\[
x_{k+1}=Ax_k,
\qquad
A=
\begin{pmatrix}
0.8&0.2\\
0.2&0.8
\end{pmatrix}.
\]
При начальном состоянии $x_0=(3,-1)^{\trans}$ найти $x_k$ и объяснить физический смысл собственных режимов.
\end{problem}

\begin{proof}[Решение]
Матрица симметрична. Для направления
\[
q_1=
\begin{pmatrix}1\\1\end{pmatrix}
\]
получаем $Aq_1=q_1$, а для
\[
q_2=
\begin{pmatrix}1\\-1\end{pmatrix}
\]
получаем $Aq_2=0.6q_2$. Первый режим изменяет обе температуры одинаково и описывает средний уровень; второй описывает разность температур.

Разложение начального состояния:
\[
\begin{pmatrix}3\\-1\end{pmatrix}
=
\begin{pmatrix}1\\1\end{pmatrix}
+2
\begin{pmatrix}1\\-1\end{pmatrix}.
\]
Поэтому
\[
x_k=
\begin{pmatrix}1\\1\end{pmatrix}
+2(0.6)^k
\begin{pmatrix}1\\-1\end{pmatrix}
=
\begin{pmatrix}
1+2(0.6)^k\\
1-2(0.6)^k
\end{pmatrix}.
\]
Средняя температура сохраняется, поскольку соответствующее собственное значение равно $1$. Температурный перепад затухает как $(0.6)^k$, поэтому обе зоны сходятся к одному отклонению, равному $1$.
\end{proof}

\begin{problem}
\textit{Практическая задача: собственные формы двух связанных узлов.}
Малые поперечные смещения двух одинаковых узлов конструкции описываются уравнением
\[
\ddot u+Ku=0,
\qquad
K=
\begin{pmatrix}
2&-1\\
-1&2
\end{pmatrix},
\]
где массы нормированы к единице. Найти собственные формы и собственные частоты.
\end{problem}

\begin{proof}[Решение]
Характеристический многочлен матрицы жёсткости:
\[
p_K(t)=
\det
\begin{pmatrix}
t-2&1\\
1&t-2
\end{pmatrix}
=(t-1)(t-3).
\]
Для $\lambda_1=1$ собственное направление задаётся вектором
\[
q_1=
\begin{pmatrix}1\\1\end{pmatrix},
\]
то есть оба узла движутся в одну сторону. Для $\lambda_2=3$ имеем
\[
q_2=
\begin{pmatrix}1\\-1\end{pmatrix},
\]
то есть узлы движутся навстречу друг другу.

Если искать движение одной моды в виде $u(t)=q\cos(\omega t)$, то подстановка даёт
\[
Kq=\omega^2q.
\]
Следовательно,
\[
\omega_1=1,
\qquad
\omega_2=\sqrt3.
\]
Противофазная форма имеет большую частоту, поскольку сильнее растягивает связь между узлами и требует большей упругой энергии.
\end{proof}

\section{Жорданова форма и функции от оператора}
\label{sec:jordan-and-functions}

Диагонализация разделяет действие оператора на независимые собственные направления, но полный базис собственных векторов существует не всегда. Жорданова форма показывает, что происходит в недиагонализируемом случае: к собственным направлениям добавляются присоединённые направления, по которым оператор последовательно передаёт состояние вдоль конечной цепочки. Эта структура объясняет появление множителей вида $k\lambda^k$ в дискретной динамике и $t e^{\lambda t}$ в решениях линейных дифференциальных уравнений.

\subsection{Почему собственных векторов может не хватить}

Рассмотрим матрицу
\[
A=
\begin{pmatrix}
\lambda&1\\
0&\lambda
\end{pmatrix}.
\]
Её характеристический многочлен равен $(t-\lambda)^2$, поэтому алгебраическая кратность собственного значения $\lambda$ равна двум. Однако
\[
A-\lambda I=
\begin{pmatrix}
0&1\\
0&0
\end{pmatrix},
\qquad
\ker(A-\lambda I)
=
\operatorname{span}
\left
\{
\begin{pmatrix}1\\0\end{pmatrix}
\right\}.
\]
Собственное подпространство одномерно. Второго независимого собственного вектора нет, поэтому матрицу нельзя диагонализировать.

Недостающее направление задаёт вектор
\[
v_2=
\begin{pmatrix}0\\1\end{pmatrix},
\qquad
(A-\lambda I)v_2=
\begin{pmatrix}1\\0\end{pmatrix}=v_1.
\]
Вектор $v_2$ не является собственным, но после одного применения $A-\lambda I$ переходит в собственный вектор $v_1$. Именно такие векторы дополняют собственные векторы до жорданова базиса.

\subsection{Корневые векторы и жордановы цепочки}

\begin{definition}
Вектор $v\in V$ называется корневым, или обобщённым собственным, вектором оператора $A$, отвечающим числу $\lambda$, если
\[
(A-\lambda I)^m v=0
\]
для некоторого натурального $m$. Наименьшее такое $m$ называется высотой вектора $v$.
\end{definition}

Собственный вектор имеет высоту один. Если высота $v$ равна $m$, то
\[
(A-\lambda I)^{m-1}v\neq0,
\qquad
(A-\lambda I)^m v=0.
\]
Последовательность
\[
v_1=(A-\lambda I)^{m-1}v,
\quad
v_2=(A-\lambda I)^{m-2}v,
\quad\dots\quad,
v_m=v
\]
удовлетворяет соотношениям
\[
(A-\lambda I)v_1=0,
\qquad
(A-\lambda I)v_{j+1}=v_j.
\]
Она называется жордановой цепочкой. Оператор действует на ней по правилу
\[
Av_1=\lambda v_1,
\qquad
Av_{j+1}=\lambda v_{j+1}+v_j.
\]
Таким образом, каждое направление масштабируется множителем $\lambda$ и одновременно передаёт часть состояния в предыдущее направление цепочки.

\begin{definition}
Корневым подпространством, отвечающим $\lambda$, называется
\[
\mathcal V_\lambda(A)
=
\bigcup_{m\geq1}\ker(A-\lambda I)^m.
\]
В конечномерном пространстве цепочка ядер стабилизируется, поэтому для некоторого $m$
\[
\mathcal V_\lambda(A)=\ker(A-\lambda I)^m.
\]
\end{definition}

Обычное собственное подпространство $V_\lambda(A)=\ker(A-\lambda I)$ содержится в $\mathcal V_\lambda(A)$. Размерность корневого подпространства равна алгебраической кратности $\lambda$, поэтому корневых векторов достаточно для заполнения той части пространства, которую нельзя заполнить обычными собственными векторами.

\subsection{Нильпотентная часть и жордановы клетки}

\begin{definition}
Оператор $N$ называется нильпотентным, если $N^r=0$ для некоторого натурального $r$. Наименьшее такое $r$ называется индексом нильпотентности.
\end{definition}

На корневом подпространстве $\mathcal V_\lambda(A)$ оператор
\[
N=A-\lambda I
\]
нильпотентен. В базисе одной жордановой цепочки длины $r$ матрица $A$ имеет вид
\[
J_r(\lambda)
=
\begin{pmatrix}
\lambda&1&0&\cdots&0\\
0&\lambda&1&\ddots&\vdots\\
\vdots&\ddots&\ddots&\ddots&0\\
0&\cdots&0&\lambda&1\\
0&\cdots&\cdots&0&\lambda
\end{pmatrix}
=
\lambda I+N_r,
\]
где $N_r$ содержит единицы над главной диагональю и удовлетворяет
\[
N_r^r=0,
\qquad
N_r^{r-1}\neq0.
\]
Матрица $J_r(\lambda)$ называется жордановой клеткой порядка $r$.

\begin{theorem}[жорданова нормальная форма]
Пусть характеристический многочлен оператора $A$ разлагается над полем $K$ на линейные множители. Тогда существует базис, в котором матрица $A$ клеточно-диагональна и имеет вид
\[
J=
\operatorname{diag}
\bigl(
J_{r_1}(\lambda_1),\dots,J_{r_s}(\lambda_s)
\bigr).
\]
Иными словами, существует обратимая матрица $S$, для которой
\[
A=SJS^{-1}.
\]
\end{theorem}

Над $\mathbb C$ условие теоремы выполняется всегда. Над $\mathbb R$ комплексные собственные значения приходится рассматривать попарно или переходить к комплексификации.

Жорданова форма хранит больше информации, чем набор собственных значений. Сумма размеров клеток со значением $\lambda$ равна его алгебраической кратности. Число таких клеток равно геометрической кратности
\[
\dim\ker(A-\lambda I).
\]
Наибольший размер клетки показывает, сколько последовательных применений $A-\lambda I$ требуется, чтобы уничтожить всё корневое подпространство.

\subsection{Минимальный многочлен}

\begin{definition}
Ненулевой многочлен $p$ называется аннулирующим для оператора $A$, если $p(A)=0$. Нормированный аннулирующий многочлен наименьшей степени называется минимальным многочленом и обозначается $m_A(t)$.
\end{definition}

Если при собственном значении $\lambda_i$ максимальный размер жордановой клетки равен $r_i$, то
\[
m_A(t)=\prod_i(t-\lambda_i)^{r_i}.
\]
Характеристический многочлен учитывает суммарные размеры всех клеток, а минимальный --- только максимальную клетку для каждого собственного значения.

\begin{criterion}[диагонализируемость через минимальный многочлен]
Оператор диагонализируем тогда и только тогда, когда его минимальный многочлен разлагается на попарно различные линейные множители.
\end{criterion}

Повторный множитель $(t-\lambda)^r$ при $r>1$ означает наличие жордановой цепочки длины больше одного. Именно она препятствует диагонализации.

\begin{theorem}[Гамильтона--Кэли]
Каждая квадратная матрица аннулируется своим характеристическим многочленом:
\[
p_A(A)=0.
\]
\end{theorem}

Теорема позволяет заменять высокие степени матрицы комбинацией меньших степеней. После деления многочлена $f(t)$ на $p_A(t)$ с остатком
\[
f(t)=q(t)p_A(t)+r(t),
\qquad
\deg r<n,
\]
получаем
\[
f(A)=r(A).
\]
На практике удобнее использовать минимальный многочлен, если он известен, поскольку его степень может быть меньше.

\subsection{Многочлены и аналитические функции от матрицы}

Для многочлена
\[
f(t)=a_0+a_1t+\cdots+a_mt^m
\]
определяется оператор
\[
f(A)=a_0I+a_1A+\cdots+a_mA^m.
\]
Это определение согласовано с заменой базиса:
\[
A=SJS^{-1}
\quad\Longrightarrow\quad
f(A)=Sf(J)S^{-1}.
\]
Поэтому достаточно уметь вычислять функцию от одной жордановой клетки.

Пусть $J=\lambda I+N$, где $N^r=0$. Для многочлена, а также для аналитической функции в окрестности $\lambda$, ряд Тейлора обрывается после степени $r-1$:
\[
f(J)
=
\sum_{s=0}^{r-1}
\frac{f^{(s)}(\lambda)}{s!}N^s.
\]
Если клетка имеет размер один, остаётся обычное правило $f(J)=f(\lambda)$. В клетке большего размера появляются производные функции. Это аналитическое отражение того, что одного значения $f(\lambda)$ недостаточно для описания действия на присоединённых направлениях.

\subsection{Степени жордановой клетки}

Для целого $k\geq0$ биномиальная формула даёт
\[
J^k=(\lambda I+N)^k
=
\sum_{s=0}^{\min(k,r-1)}
\binom{k}{s}\lambda^{k-s}N^s.
\]
Для клетки второго порядка
\[
J_2(\lambda)^k
=
\begin{pmatrix}
\lambda^k&k\lambda^{k-1}\\
0&\lambda^k
\end{pmatrix}.
\]
Поэтому повторное применение недиагонализируемой матрицы создаёт не только геометрическую прогрессию $\lambda^k$, но и полиномиальные множители по $k$.

При $|\lambda|<1$ экспоненциальное затухание в итоге сильнее любого полиномиального множителя, и клетка стремится к нулю. Однако на первых шагах множитель $k$ может вызвать заметный переходный рост. Если $|\lambda|=1$ и размер клетки больше единицы, степени матрицы могут расти полиномиально, несмотря на то что модуль собственного значения равен единице.

\subsection{Матричная экспонента и линейная динамика}

\begin{definition}
Экспонентой матрицы $A$ называется абсолютно сходящийся ряд
\[
e^A
=
I+A+\frac{A^2}{2!}+\frac{A^3}{3!}+\cdots.
\]
\end{definition}

Для жордановой клетки
\[
e^{tJ}
=
e^{\lambda t}
\sum_{s=0}^{r-1}\frac{t^s}{s!}N^s.
\]
В частности,
\[
e^{tJ_2(\lambda)}
=
e^{\lambda t}
\begin{pmatrix}
1&t\\
0&1
\end{pmatrix}.
\]

Матричная экспонента является фундаментальным решением системы
\[
x'(t)=Ax(t),
\qquad
x(0)=x_0:
\]
действительно,
\[
x(t)=e^{tA}x_0.
\]
Если $A=SJS^{-1}$, то
\[
e^{tA}=Se^{tJ}S^{-1}.
\]
Таким образом, жордановы клетки объясняют, почему решения линейной системы состоят из произведений экспонент $e^{\lambda t}$ на многочлены по $t$.

\begin{computationalnote}
Жорданова форма важна для точного анализа, но плохо устойчива к ошибкам входных данных: малая погрешность может изменить число и размеры клеток. В вычислениях с числами с плавающей точкой обычно используют ортогональные или унитарные разложения, прежде всего форму Шура. Жорданова форма здесь служит языком понимания структуры, а не основным численным алгоритмом.
\end{computationalnote}

\subsection*{Задачи}

\begin{problem}
\textit{Задача Винберга.}
Доказать, что высота любого корневого вектора, отвечающего собственному значению $\lambda$, не превосходит $\dim\mathcal V_\lambda(A)$.
\end{problem}

\begin{proof}[Решение]
Пусть корневой вектор $v$ имеет высоту $m$. Тогда
\[
(A-\lambda I)^m v=0,
\qquad
(A-\lambda I)^{m-1}v\neq0.
\]
Обозначим $N=A-\lambda I$. Покажем, что векторы
\[
v,Nv,N^2v,\dots,N^{m-1}v
\]
линейно независимы. Пусть
\[
a_0v+a_1Nv+\cdots+a_{m-1}N^{m-1}v=0
\]
и $a_j$ --- первый ненулевой коэффициент. Применение $N^{m-1-j}$ уничтожает все слагаемые после $a_jN^jv$ и даёт
\[
a_jN^{m-1}v=0,
\]
что невозможно. Следовательно, имеется $m$ независимых векторов, и все они лежат в $\mathcal V_\lambda(A)$. Поэтому
\[
m\leq\dim\mathcal V_\lambda(A).
\]
\end{proof}

\begin{problem}
\textit{Задача Винберга.}
Доказать, что максимальный порядок жордановых клеток с собственным значением $\lambda$ равен индексу нильпотентности оператора
\[
N=(A-\lambda I)|_{\mathcal V_\lambda(A)}.
\]
\end{problem}

\begin{proof}[Решение]
В жордановом базисе оператор $N$ является клеточно-диагональным и состоит из нильпотентных клеток
\[
N_{r_1},\dots,N_{r_s}.
\]
Для клетки порядка $r_i$ выполнено
\[
N_{r_i}^{r_i}=0,
\qquad
N_{r_i}^{r_i-1}\neq0.
\]
Степень всей клеточно-диагональной матрицы равна нулю тогда и только тогда, когда обнулилась каждая клетка. Поэтому
\[
N^m=0
\quad\Longleftrightarrow\quad
m\geq\max_i r_i.
\]
Наименьшее такое $m$ равно $\max_i r_i$, то есть максимальному размеру жордановой клетки.
\end{proof}

\begin{problem}
\textit{Практическая задача: привод с внутренним запаздыванием.}
Пусть $u(t)$ --- остаточное управляющее воздействие привода, а $y(t)$ --- отклонение положения панели. Линеаризованная модель имеет вид
\[
\begin{pmatrix}y\\u\end{pmatrix}'
=
A
\begin{pmatrix}y\\u\end{pmatrix},
\qquad
A=
\begin{pmatrix}
-1&1\\
0&-1
\end{pmatrix}.
\]
Найти движение при $y(0)=0$, $u(0)=1$ и объяснить появление множителя $t$.
\end{problem}

\begin{proof}[Решение]
Матрица является жордановой клеткой
\[
A=-I+N,
\qquad
N=
\begin{pmatrix}
0&1\\
0&0
\end{pmatrix},
\qquad
N^2=0.
\]
Поэтому
\[
e^{tA}=e^{-t}(I+tN)
=e^{-t}
\begin{pmatrix}
1&t\\
0&1
\end{pmatrix}.
\]
Для начального состояния $x_0=(0,1)^{\trans}$ получаем
\[
\begin{pmatrix}y(t)\\u(t)\end{pmatrix}
=
e^{tA}x_0
=
\begin{pmatrix}
t e^{-t}\\e^{-t}
\end{pmatrix}.
\]
Внутреннее воздействие $u$ затухает как $e^{-t}$, но некоторое время продолжает передаваться в координату $y$. Интегрирование этого затухающего сигнала создаёт множитель $t$. Поэтому положение сначала отклоняется, а затем возвращается к нулю, хотя оба собственных значения равны $-1$.
\end{proof}

\begin{problem}
\textit{Практическая задача: регулятор с памятью.}
Пусть $e_k$ --- ошибка формы на $k$-й итерации, а $m_k$ --- внутреннее состояние памяти регулятора. Обновление задаётся системой
\[
\begin{pmatrix}e_{k+1}\\m_{k+1}\end{pmatrix}
=
A
\begin{pmatrix}e_k\\m_k\end{pmatrix},
\qquad
A=
\begin{pmatrix}
0.9&1\\
0&0.9
\end{pmatrix}.
\]
Найти состояние при $(e_0,m_0)=(0,1)$ и объяснить, почему ошибка может сначала вырасти, хотя единственное собственное значение по модулю меньше единицы.
\end{problem}

\begin{proof}[Решение]
Матрица является клеткой $J_2(0.9)$. Для $k\geq1$
\[
A^k
=
\begin{pmatrix}
0.9^k&k\,0.9^{k-1}\\
0&0.9^k
\end{pmatrix}.
\]
Следовательно,
\[
\begin{pmatrix}e_k\\m_k\end{pmatrix}
=A^k
\begin{pmatrix}0\\1\end{pmatrix}
=
\begin{pmatrix}
k\,0.9^{k-1}\\0.9^k
\end{pmatrix}.
\]
Память $m_k$ монотонно затухает, но на каждом шаге добавляется к ошибке формы. Поэтому $e_k$ содержит полиномиальный множитель $k$ и сначала может увеличиваться. В пределе экспоненциальное затухание $0.9^k$ побеждает линейный рост, так что $e_k\to0$. Один только спектральный радиус сообщает предельное поведение, но не описывает переходный рост недиагонализируемой системы.
\end{proof}

\section{Аффинные пространства и выпуклость}
\label{sec:affine-convex}

Векторное пространство содержит выделенный нулевой вектор. Для пространства точек такая привилегированная точка обычно отсутствует: положение объекта зависит от выбора начала координат, а вектор перемещения --- нет. Аффинная геометрия отделяет точки от векторов и сохраняет те конструкции, которые не зависят от произвольного выбора начала отсчёта.

\subsection{Точки, векторы и выбор начала отсчёта}

\begin{definition}
Аффинным пространством, ассоциированным с векторным пространством \(V\), называется множество точек \(S\), для которого определена операция
\[
S\times V\to S,
\qquad
(p,v)\mapsto p+v,
\]
такая, что
\[
(p+u)+v=p+(u+v)
\]
и для любых \(p,q\in S\) существует единственный вектор \(v\in V\), удовлетворяющий равенству
\[
q=p+v.
\]
Этот вектор обозначается через
\[
\overrightarrow{pq}=q-p.
\]
\end{definition}

Точки нельзя складывать между собой как векторы без дополнительных условий. Корректны разность двух точек, сумма точки и вектора, а также комбинации точек, коэффициенты которых в сумме дают единицу.

Выбор точки \(o\in S\) в качестве начала отсчёта позволяет отождествить \(S\) с \(V\):
\[
p\longmapsto \overrightarrow{op}.
\]
Другой выбор начала изменит координаты всех точек на один и тот же перенос, но не изменит разности \(\overrightarrow{pq}\).

\begin{example}
Координата точки на прямой зависит от положения нуля. Расстояние и направленное смещение между двумя точками от этого выбора не зависят. Именно поэтому положение является аффинной величиной, а перемещение --- векторной.
\end{example}

\subsection{Аффинные подпространства и аффинная оболочка}

\begin{definition}
Пусть \(U\subseteq V\) --- линейное подпространство и \(p_0\in S\). Множество
\[
P=p_0+U
=
\{p_0+u\mid u\in U\}
\]
называется аффинным подпространством с направляющим пространством \(U\).
\end{definition}

Если \(0\notin P\), то \(P\) не является линейным подпространством после выбора произвольных координат. Однако все векторы между его точками принадлежат одному и тому же пространству \(U\).

\begin{definition}
Аффинной оболочкой множества \(M\subseteq S\) называется наименьшее аффинное подпространство, содержащее \(M\). Она обозначается через
\[
\operatorname{aff}M.
\]
\end{definition}

Для непустого \(M\) и любой фиксированной точки \(p_0\in M\)
\[
\operatorname{aff}M
=
p_0+
\operatorname{span}
\left\{
\overrightarrow{p_0p}\mid p\in M
\right\}.
\]
Эта формула переводит аффинную задачу в линейную: выбирается одна опорная точка, после чего исследуются векторы от неё к остальным точкам.

\begin{proposition}
Через \(k+1\) точек проходит аффинное подпространство размерности не более \(k\). Его размерность равна \(k\) тогда и только тогда, когда векторы
\[
\overrightarrow{p_0p_1},\dots,\overrightarrow{p_0p_k}
\]
линейно независимы.
\end{proposition}

\subsection{Аффинная независимость и барицентрические координаты}

\begin{definition}
Точки \(p_0,\dots,p_k\) называются аффинно независимыми, если векторы
\[
\overrightarrow{p_0p_1},\dots,\overrightarrow{p_0p_k}
\]
линейно независимы.
\end{definition}

Две различные точки аффинно независимы и задают прямую. Три точки аффинно независимы тогда и только тогда, когда не лежат на одной прямой. Четыре точки в трёхмерном пространстве аффинно независимы тогда и только тогда, когда не лежат в одной плоскости.

\begin{definition}
Выражение
\[
p=\lambda_0p_0+\cdots+\lambda_kp_k,
\qquad
\lambda_0+\cdots+\lambda_k=1,
\]
называется аффинной, или барицентрической, комбинацией точек.
\end{definition}

Смысл формальной записи определяется выбором любой опорной точки, например \(p_0\):
\[
p
=
p_0+
\sum_{i=1}^{k}\lambda_i\overrightarrow{p_0p_i},
\qquad
\lambda_0=1-\sum_{i=1}^{k}\lambda_i.
\]
Условие суммы коэффициентов, равной единице, делает результат независимым от выбора начала координат.

\begin{theorem}
Если точки \(p_0,\dots,p_k\) аффинно независимы, то каждая точка их аффинной оболочки имеет единственные барицентрические координаты:
\[
p=\sum_{i=0}^{k}\lambda_ip_i,
\qquad
\sum_{i=0}^{k}\lambda_i=1.
\]
\end{theorem}

Барицентрические координаты показывают положение точки относительно вершин симплекса. Они применяются в конечных элементах, интерполяции полей, геометрическом моделировании и представлении смесей.

\subsection{Аффинные отображения}

\begin{definition}
Отображение \(f:S\to S'\) называется аффинным, если существует линейное отображение \(L:V\to V'\), для которого
\[
f(p+v)=f(p)+L(v)
\]
при всех \(p\in S\) и \(v\in V\). Линейное отображение \(L\) называется линейной частью, или дифференциалом, аффинного отображения.
\end{definition}

После выбора начал координат аффинное отображение записывается в привычном виде
\[
f(x)=Ax+b.
\]
Матрица \(A\) описывает изменение направлений, а вектор \(b\) --- общий перенос. Отображение обратимо тогда и только тогда, когда обратима его линейная часть \(A\).

\begin{proposition}
Аффинное отображение сохраняет барицентрические комбинации:
\[
f\left(\sum_{i=0}^{k}\lambda_ip_i\right)
=
\sum_{i=0}^{k}\lambda_if(p_i),
\qquad
\sum_{i=0}^{k}\lambda_i=1.
\]
\end{proposition}

Именно это свойство делает возможным перенос интерполяции с эталонного геометрического элемента на реальный. Достаточно знать образы вершин, после чего любая внутренняя точка переносится с теми же барицентрическими коэффициентами.

\subsection{Выпуклые комбинации и выпуклые множества}

\begin{definition}
Барицентрическая комбинация
\[
p=\sum_{i=0}^{k}\lambda_ip_i
\]
называется выпуклой, если
\[
\lambda_i\geq0,
\qquad
\sum_{i=0}^{k}\lambda_i=1.
\]
\end{definition}

Для двух точек выпуклые комбинации имеют вид
\[
(1-t)p+tq,
\qquad
0\leq t\leq1,
\]
и заполняют весь отрезок между ними.

\begin{definition}
Множество \(C\subseteq S\) называется выпуклым, если вместе с любыми двумя точками \(p,q\in C\) оно содержит весь соединяющий их отрезок:
\[
(1-t)p+tq\in C
\quad
\text{для всех }t\in[0,1].
\]
\end{definition}

\begin{proposition}
Выпуклое множество содержит любую конечную выпуклую комбинацию своих точек.
\end{proposition}

\begin{definition}
Выпуклой оболочкой множества \(M\) называется множество всех конечных выпуклых комбинаций его точек:
\[
\operatorname{conv}M
=
\left\{
\sum_{i=1}^{m}\lambda_ip_i
\;\middle|\;
p_i\in M,\
\lambda_i\geq0,\
\sum_{i=1}^{m}\lambda_i=1
\right\}.
\]
\end{definition}

Выпуклая оболочка является наименьшим выпуклым множеством, содержащим \(M\). В отличие от аффинной оболочки она запрещает отрицательные коэффициенты, поэтому не продолжает комбинации за границы, заданные исходными точками.

\subsection{Симплексы и проверка принадлежности}

\begin{definition}
Выпуклая оболочка \(k+1\) аффинно независимых точек
\[
\Delta=\operatorname{conv}\{p_0,\dots,p_k\}
\]
называется \(k\)-мерным симплексом.
\end{definition}

Отрезок является одномерным симплексом, треугольник вместе с внутренностью --- двумерным, тетраэдр вместе с внутренностью --- трёхмерным.

\begin{criterion}[барицентрический критерий]
Пусть \(p_0,\dots,p_k\) аффинно независимы, а
\[
p=\sum_{i=0}^{k}\lambda_ip_i,
\qquad
\sum_{i=0}^{k}\lambda_i=1.
\]
Тогда
\[
p\in\operatorname{conv}\{p_0,\dots,p_k\}
\]
тогда и только тогда, когда все \(\lambda_i\geq0\).
\end{criterion}

Если один коэффициент отрицателен, точка лежит за гранью, противоположной соответствующей вершине. Если некоторый коэффициент равен нулю, точка принадлежит соответствующей грани симплекса.

\subsection{Полупространства, многогранники и крайние точки}

Пусть \(a\in\R^n\), \(a\neq0\), и \(\beta\in\R\). Множество
\[
H^-=\{x\in\R^n\mid a^{\trans}x\leq\beta\}
\]
является полупространством. Оно выпукло, поскольку линейное неравенство сохраняется при выпуклом смешивании допустимых точек.

\begin{definition}
Пересечение конечного числа полупространств называется выпуклым многогранником:
\[
P=\{x\in\R^n\mid Ax\leq b\}.
\]
\end{definition}

Так записываются линейные ограничения на размеры, массу, ресурсы и управляющие воздействия. Выпуклость означает, что смесь двух допустимых проектов остаётся допустимой в рамках линейной модели.

\begin{definition}
Точка \(p\in C\) называется крайней, если из представления
\[
p=(1-t)x+ty,
\qquad
x,y\in C,\quad 0<t<1,
\]
следует \(x=y=p\).
\end{definition}

Крайние точки многогранника являются его вершинами. Если линейная функция достигает максимума на ограниченном многограннике, то хотя бы одна вершина также является точкой максимума. На этом основана геометрическая идея линейного программирования: оптимальное решение можно искать среди граничных конфигураций, а не во всём непрерывном множестве.

\paragraph{Связь с генеративным проектированием.}

Аффинные координаты позволяют отделить геометрию объекта от способа её записи. Барицентрическая интерполяция переносит значения между вершинами сетки и внутренними точками. Выпуклая оболочка описывает все смеси прототипов с неотрицательными весами, а линейные неравенства формируют приближённую область допустимых конструкций.

При этом выпуклость является свойством математической модели, а не автоматическим свойством физической системы. Линейная смесь параметров может не соответствовать линейной смеси прочности, устойчивости или технологичности. Поэтому выпуклая модель полезна как вычислительно прозрачное приближение, но её физические предположения должны проверяться отдельно.

\subsection*{Задачи}

\begin{problem}
\textit{Задача по материалу Винберга.}
Пусть \(p_0,\dots,p_k\) --- аффинно независимые точки. Доказать единственность их барицентрических координат.
\end{problem}

\begin{proof}[Решение]
Предположим, что одна точка имеет два представления:
\[
p=\sum_{i=0}^{k}\lambda_ip_i
=
\sum_{i=0}^{k}\mu_ip_i,
\qquad
\sum_{i=0}^{k}\lambda_i
=
\sum_{i=0}^{k}\mu_i
=1.
\]
Вычтем представления и используем \(p_0\) как опорную точку:
\[
0
=
\sum_{i=0}^{k}(\lambda_i-\mu_i)p_i
=
\sum_{i=1}^{k}(\lambda_i-\mu_i)
\overrightarrow{p_0p_i}.
\]
Векторы
\[
\overrightarrow{p_0p_1},\dots,\overrightarrow{p_0p_k}
\]
линейно независимы, поэтому
\[
\lambda_i-\mu_i=0
\quad
\text{при }i=1,\dots,k.
\]
Из условия суммы коэффициентов следует также
\[
\lambda_0=\mu_0.
\]
Следовательно, все барицентрические координаты совпадают.
\end{proof}

\begin{problem}
\textit{Задача по материалу Винберга.}
Пусть \(f:S\to S'\) --- аффинное отображение. Доказать, что для любого множества \(M\subseteq S\)
\[
f(\operatorname{conv}M)
=
\operatorname{conv}f(M).
\]
\end{problem}

\begin{proof}[Решение]
Возьмём \(p\in\operatorname{conv}M\). Тогда
\[
p=\sum_{i=1}^{m}\lambda_ip_i,
\qquad
p_i\in M,\quad
\lambda_i\geq0,\quad
\sum_{i=1}^{m}\lambda_i=1.
\]
Аффинность даёт
\[
f(p)
=
\sum_{i=1}^{m}\lambda_if(p_i),
\]
поэтому \(f(p)\in\operatorname{conv}f(M)\). Получено включение
\[
f(\operatorname{conv}M)
\subseteq
\operatorname{conv}f(M).
\]

Обратно, любая точка \(q\in\operatorname{conv}f(M)\) представляется как
\[
q=\sum_{i=1}^{m}\lambda_if(p_i)
=
f\left(\sum_{i=1}^{m}\lambda_ip_i\right),
\]
причём внутренняя сумма принадлежит \(\operatorname{conv}M\). Значит,
\[
\operatorname{conv}f(M)
\subseteq
f(\operatorname{conv}M).
\]
Оба включения дают требуемое равенство.
\end{proof}

\begin{problem}
\textit{Практическая задача: интерполяция на конечном элементе.}
Эталонный треугольник имеет вершины
\[
\widehat p_0=(0,0),
\qquad
\widehat p_1=(1,0),
\qquad
\widehat p_2=(0,1).
\]
Он аффинно отображается в физический элемент с вершинами
\[
p_0=(1,1),
\qquad
p_1=(5,1),
\qquad
p_2=(1,3).
\]
Найти образ точки
\[
\widehat p=(0.25,0.5)
\]
и интерполированную температуру, если в вершинах физического элемента заданы значения
\[
T_0=20,\qquad T_1=28,\qquad T_2=24.
\]
\end{problem}

\begin{proof}[Решение]
В эталонном треугольнике
\[
\widehat p
=
\lambda_0\widehat p_0
+\lambda_1\widehat p_1
+\lambda_2\widehat p_2.
\]
Из координат сразу получаем
\[
\lambda_1=0.25,
\qquad
\lambda_2=0.5,
\qquad
\lambda_0=1-0.25-0.5=0.25.
\]
Все коэффициенты неотрицательны, поэтому точка лежит внутри треугольника.

Аффинное отображение сохраняет барицентрические координаты:
\[
p
=
0.25p_0+0.25p_1+0.5p_2.
\]
Следовательно,
\[
p
=
0.25(1,1)+0.25(5,1)+0.5(1,3)
=
(2,2).
\]
Линейная интерполяция температуры использует те же коэффициенты:
\[
T(p)
=
0.25\cdot20
+0.25\cdot28
+0.5\cdot24
=
24.
\]
Одна система барицентрических координат одновременно переносит геометрию и поле температуры.
\end{proof}

\begin{problem}
\textit{Практическая задача: смесь трёх вариантов материала.}
Упрощённая модель материала описывает нормированные жёсткость и плотность парой
\[
(E,\rho).
\]
Доступны три базовых варианта:
\[
A=(1,1),
\qquad
B=(3,1),
\qquad
C=(1,2).
\]
Можно ли получить целевые характеристики
\[
P=(2,1.4)
\]
в виде выпуклой смеси трёх вариантов? Найти доли смеси.
\end{problem}

\begin{proof}[Решение]
Ищем коэффициенты
\[
P=\lambda_AA+\lambda_BB+\lambda_CC,
\qquad
\lambda_A+\lambda_B+\lambda_C=1,
\qquad
\lambda_A,\lambda_B,\lambda_C\geq0.
\]
По координате жёсткости
\[
\lambda_A+3\lambda_B+\lambda_C=2.
\]
Вычитая условие суммы коэффициентов, получаем
\[
2\lambda_B=1,
\qquad
\lambda_B=0.5.
\]
По координате плотности
\[
\lambda_A+\lambda_B+2\lambda_C=1.4.
\]
Снова вычитаем сумму коэффициентов:
\[
\lambda_C=0.4.
\]
Тогда
\[
\lambda_A=1-0.5-0.4=0.1.
\]
Все коэффициенты неотрицательны, поэтому целевая точка принадлежит треугольнику
\[
\operatorname{conv}\{A,B,C\}.
\]
В рамках линейной модели требуются доли
\[
10\% \text{ материала }A,\qquad
50\% \text{ материала }B,\qquad
40\% \text{ материала }C.
\]
Результат проверяет достижимость целевых характеристик только в принятой линейной модели смешивания; физическую реализуемость такой смеси необходимо подтверждать отдельно.
\end{proof}

\section{Тензорная алгебра}
\label{sec:tensor-algebra}

Матрица описывает линейную зависимость между двумя наборами координат. Во многих задачах этого недостаточно: результат зависит сразу от нескольких независимых аргументов, а данные имеют несколько смысловых осей. Тензорный язык позволяет работать с такими объектами без сведения всех индексов к одному длинному вектору.

Главная идея состоит не в том, что тензор является многомерным массивом. Массив содержит координаты тензора в выбранных базисах. Сам тензор определяется тем, как он преобразуется при смене координат и как действует на векторы и линейные функционалы.

\subsection{Полилинейные отображения}

\begin{definition}
Пусть \(V_1,\dots,V_p,U\) --- векторные пространства над одним полем \(K\). Отображение
\[
\Phi:V_1\times\cdots\times V_p\to U
\]
называется \(p\)-линейным, или полилинейным, если оно линейно по каждому аргументу при фиксированных остальных аргументах.
\end{definition}

Для билинейного отображения \(\Phi:V\times W\to U\) это означает
\[
\Phi(\alpha v_1+\beta v_2,w)
=
\alpha\Phi(v_1,w)+\beta\Phi(v_2,w)
\]
и аналогичное равенство по второму аргументу.

Скалярное произведение, билинейная форма, матричное умножение и определитель являются полилинейными конструкциями. Полилинейность позволяет восстановить значение отображения по его значениям на наборах базисных векторов.

Если
\[
\dim V_i=n_i,
\qquad
\dim U=m,
\]
то пространство всех \(p\)-линейных отображений \(V_1\times\cdots\times V_p\to U\) имеет размерность
\[
n_1\cdots n_p\,m.
\]

\subsection{Тензорное произведение}

\begin{definition}
Тензорным произведением пространств \(V\) и \(W\) называется векторное пространство \(V\otimes W\) вместе с билинейным отображением
\[
V\times W\to V\otimes W,
\qquad
(v,w)\mapsto v\otimes w,
\]
обладающим следующим универсальным свойством: для любого билинейного отображения
\[
\Phi:V\times W\to U
\]
существует единственное линейное отображение
\[
\widetilde{\Phi}:V\otimes W\to U,
\]
для которого
\[
\Phi(v,w)=\widetilde{\Phi}(v\otimes w).
\]
\end{definition}

Универсальное свойство переводит билинейную задачу в линейную. Вместо двух аргументов \(v,w\) рассматривается один объект \(v\otimes w\), после чего применяется обычное линейное отображение.

Из билинейности следуют соотношения
\[
(v_1+v_2)\otimes w
=
v_1\otimes w+v_2\otimes w,
\]
\[
v\otimes(w_1+w_2)
=
v\otimes w_1+v\otimes w_2,
\]
\[
(\lambda v)\otimes w
=
v\otimes(\lambda w)
=
\lambda(v\otimes w).
\]
При этом обычно
\[
(v_1+v_2)\otimes(w_1+w_2)
\neq
v_1\otimes w_1+v_2\otimes w_2,
\]
поскольку в раскрытии присутствуют два смешанных слагаемых.

\begin{theorem}
Если \(e_1,\dots,e_n\) --- базис \(V\), а \(f_1,\dots,f_m\) --- базис \(W\), то
\[
e_i\otimes f_j,
\qquad
1\leq i\leq n,\quad 1\leq j\leq m,
\]
образуют базис пространства \(V\otimes W\). Следовательно,
\[
\dim(V\otimes W)=\dim V\cdot\dim W.
\]
\end{theorem}

Каждый тензор \(T\in V\otimes W\) единственным образом записывается как
\[
T=\sum_{i=1}^{n}\sum_{j=1}^{m}t_{ij}\,e_i\otimes f_j.
\]
Коэффициенты \(t_{ij}\) удобно располагать в прямоугольной таблице. Поэтому тензор второго порядка после выбора базисов выглядит как матрица.

\subsection{Разложимые тензоры и тензорный ранг}

Тензор вида
\[
v\otimes w
\]
называется разложимым, или простым. Если
\[
v=\sum_i v_i e_i,
\qquad
w=\sum_j w_j f_j,
\]
то
\[
v\otimes w
=
\sum_{i,j}v_iw_j\,e_i\otimes f_j.
\]
Его координатная матрица равна диадному произведению координатных столбцов:
\[
[v\otimes w]=[v][w]^{\trans}.
\]
Здесь используется термин \emph{диадное произведение}. Он не относится к внешнему произведению \(u\wedge v\), которое будет введено ниже в алгебре Грассмана. Такая матрица имеет ранг не выше единицы.

Не каждый тензор является разложимым. В общем случае
\[
T=\sum_{r=1}^{R}v_r\otimes w_r.
\]
Наименьшее возможное число слагаемых \(R\) называется тензорным рангом для тензора второго порядка и совпадает с матричным рангом его координатной матрицы. Для тензоров более высокого порядка понятие ранга становится сложнее и перестаёт сводиться к обычному рангу матрицы.

\subsection{Тензорное произведение линейных отображений}

Пусть
\[
A:V\to V',
\qquad
B:W\to W'
\]
--- линейные отображения. Их тензорное произведение определяется правилом
\[
(A\otimes B)(v\otimes w)
=
Av\otimes Bw
\]
и продолжается по линейности на всё \(V\otimes W\).

Если матрицы \(A\) и \(B\) выбраны в соответствующих базисах, то матрица \(A\otimes B\) имеет блочный вид
\[
A\otimes B
=
\begin{pmatrix}
a_{11}B&\cdots&a_{1n}B\\
\vdots&\ddots&\vdots\\
a_{m1}B&\cdots&a_{mn}B
\end{pmatrix}.
\]
Эта матрица называется произведением Кронекера. Она возникает при построении операторов на прямоугольных сетках, раздельной обработке пространственных и канальных координат и описании составных систем.

\subsection{Тензоры типа \texorpdfstring{\((p,q)\)}{(p,q)}}

Для конечномерного пространства \(V\) тензором типа \((p,q)\) называется элемент
\[
T\in
\underbrace{V\otimes\cdots\otimes V}_{p\text{ раз}}
\otimes
\underbrace{V^*\otimes\cdots\otimes V^*}_{q\text{ раз}}.
\]
Верхние индексы относятся к копиям \(V\), нижние --- к копиям \(V^*\). В базисе тензор задаётся координатами
\[
T^{i_1\dots i_p}_{j_1\dots j_q}.
\]

Вектор имеет тип \((1,0)\), линейный функционал --- тип \((0,1)\), билинейная форма --- тип \((0,2)\), линейный оператор --- тип \((1,1)\).

При смене базиса верхние и нижние индексы преобразуются противоположным образом. Это обеспечивает независимость самого тензора от координатной записи.

\subsection{Свёртка и работа с индексами}

Если у тензора имеется хотя бы один верхний и один нижний индекс, их можно свернуть:
\[
S^{i_2\dots i_p}_{j_2\dots j_q}
=
\sum_k
T^{k i_2\dots i_p}_{k j_2\dots j_q}.
\]
Свёртка уменьшает тип тензора с \((p,q)\) до \((p-1,q-1)\).

След матрицы является свёрткой тензора типа \((1,1)\):
\[
\operatorname{tr}A=A^i_i.
\]
Умножение матриц также можно читать как тензорное произведение с последующей свёрткой:
\[
(AB)^i_j
=
A^i_kB^k_j.
\]
Повторяющийся один раз сверху и один раз снизу индекс суммируется. Эта договорённость называется соглашением Эйнштейна.

В евклидовом пространстве скалярное произведение задаёт тензор
\[
g_{ij}.
\]
Он позволяет преобразовывать векторы в функционалы:
\[
v_i=g_{ij}v^j.
\]
Обратная матрица \(g^{ij}\) выполняет обратную операцию:
\[
v^i=g^{ij}v_j.
\]
В ортонормированном базисе \(g_{ij}=\delta_{ij}\), поэтому численно верхние и нижние координаты совпадают. Это удобство связано именно с выбранным базисом, а не с исчезновением различия между векторами и функционалами.

\subsection{Тензорная алгебра}

Все степени
\[
T^0(V)=K,\qquad
T^1(V)=V,\qquad
T^2(V)=V\otimes V,\quad\dots
\]
объединяются в прямую сумму
\[
T(V)=
\bigoplus_{p=0}^{\infty}T^p(V).
\]
Произведением тензоров служит тензорное произведение:
\[
T^p(V)\times T^q(V)\to T^{p+q}(V).
\]
В общем случае оно некоммутативно:
\[
u\otimes v\neq v\otimes u.
\]

Тензорная алгебра содержит одновременно все конечные последовательности векторных аргументов. Симметрическая и внешняя алгебры получаются из неё введением дополнительных соотношений.

\subsection{Симметрические тензоры}

\begin{definition}
Тензор \(T\in T^p(V)\) называется симметрическим, если его значение не меняется при любой перестановке аргументов.
\end{definition}

Симметрическое произведение обозначается через
\[
v_1\vee\cdots\vee v_p.
\]
Для двух векторов при характеристике поля, отличной от \(2\),
\[
u\vee v
=
\frac12(u\otimes v+v\otimes u).
\]

Пространство симметрических тензоров степени \(p\) обозначается через
\[
S^p(V).
\]
Если \(\dim V=n\), то
\[
\dim S^p(V)
=
\binom{n+p-1}{p}.
\]

Симметрические тензоры естественно возникают как коэффициенты однородных многочленов и производные высших порядков. Например, гессиан гладкой скалярной функции является симметрическим тензором второго порядка.

\subsection{Внешняя алгебра}

\begin{definition}
Внешним произведением векторов \(v_1,\dots,v_p\) называется антисимметрический тензор
\[
v_1\wedge\cdots\wedge v_p,
\]
который меняет знак при перестановке двух аргументов.
\end{definition}

Из антисимметрии следует
\[
v\wedge v=0,
\qquad
u\wedge v=-v\wedge u.
\]
Пространство внешних \(p\)-векторов обозначается через
\[
\Lambda^p(V).
\]

Если \(e_1,\dots,e_n\) --- базис \(V\), то базис \(\Lambda^p(V)\) образуют элементы
\[
e_{i_1}\wedge\cdots\wedge e_{i_p},
\qquad
i_1<\cdots<i_p.
\]
Поэтому
\[
\dim\Lambda^p(V)=\binom{n}{p}.
\]

\begin{theorem}
Для векторов \(v_1,\dots,v_p\)
\[
v_1\wedge\cdots\wedge v_p=0
\]
тогда и только тогда, когда система \(v_1,\dots,v_p\) линейно зависима.
\end{theorem}

В двумерном пространстве коэффициент \(u\wedge v\) в базисе \(e_1\wedge e_2\) равен определителю матрицы координат \(u,v\), то есть ориентированной площади параллелограмма. В трёхмерном пространстве \(u\wedge v\wedge w\) кодирует ориентированный объём.

Если строки матрицы являются координатами \(p\) векторов, то координаты их внешнего произведения равны минорам порядка \(p\). Эти координаты называют координатами Плюккера; они описывают само подпространство, не завися от выбора его базиса с точностью до общего ненулевого множителя.

\paragraph{Практический смысл тензорного языка.}

Тензорная запись нужна тогда, когда оси данных имеют различный смысл и преобразуются независимо. Пространственная координата, номер канала, временной индекс и тип физической величины не обязаны объединяться в один неструктурированный индекс.

Однако само наличие многомерного массива ещё не делает вычисление содержательно тензорным. Важно указать пространства, типы индексов и правила преобразования. Без этого слово ``тензор'' описывает только форму хранения данных.

\subsection*{Задачи}

\begin{problem}
\textit{Задача по материалу Винберга.}
Пусть \(V\) имеет базис \(e_1,e_2\), а \(W\) --- базис \(f_1,f_2,f_3\). Построить базис \(V\otimes W\), найти его размерность и координаты тензора
\[
T=(2e_1-e_2)\otimes(f_1+3f_3).
\]
Определить ранг координатной матрицы \(T\).
\end{problem}

\begin{proof}[Решение]
Базис тензорного произведения состоит из шести элементов:
\[
e_1\otimes f_1,\quad
e_1\otimes f_2,\quad
e_1\otimes f_3,\quad
e_2\otimes f_1,\quad
e_2\otimes f_2,\quad
e_2\otimes f_3.
\]
Следовательно,
\[
\dim(V\otimes W)=2\cdot3=6.
\]

Раскрывая произведение по билинейности, получаем
\[
T
=
2e_1\otimes f_1
+6e_1\otimes f_3
-e_2\otimes f_1
-3e_2\otimes f_3.
\]
Его координатная матрица равна
\[
[T]
=
\begin{pmatrix}
2&0&6\\
-1&0&-3
\end{pmatrix}
=
\begin{pmatrix}
2\\-1
\end{pmatrix}
\begin{pmatrix}
1&0&3
\end{pmatrix}.
\]
Вторая строка пропорциональна первой, поэтому
\[
\operatorname{rank}[T]=1.
\]
Это соответствует тому, что \(T\) является одним разложимым тензором.
\end{proof}

\begin{problem}
\textit{Задача по материалу Винберга.}
Для векторов
\[
u=(1,2,0),
\qquad
v=(0,1,1),
\qquad
w=(1,3,1)
\]
вычислить \(u\wedge v\wedge w\) и определить, линейно ли независима система.
\end{problem}

\begin{proof}[Решение]
В базисе \(e_1\wedge e_2\wedge e_3\) коэффициент внешнего произведения равен определителю:
\[
u\wedge v\wedge w
=
\det
\begin{pmatrix}
1&2&0\\
0&1&1\\
1&3&1
\end{pmatrix}
e_1\wedge e_2\wedge e_3.
\]
Вычисляем:
\[
\det
\begin{pmatrix}
1&2&0\\
0&1&1\\
1&3&1
\end{pmatrix}
=
1(1-3)-2(0-1)=0.
\]
Следовательно,
\[
u\wedge v\wedge w=0.
\]
Система линейно зависима. Действительно,
\[
w=u+v.
\]
Внешнее произведение обнаружило вырождение ориентированного объёма, не требуя отдельного решения системы уравнений.
\end{proof}

\begin{problem}
\textit{Практическая задача: пространственно-канальный признак.}
В двух пространственных областях детали измеряется трёхкомпонентный состав материала. Интенсивность областей задаётся вектором
\[
s=
\begin{pmatrix}
2\\1
\end{pmatrix},
\]
а относительный состав компонентов ---
\[
c=
\begin{pmatrix}
0.5\\0.3\\0.2
\end{pmatrix}.
\]
Построить разложимый тензор \(F=s\otimes c\), объяснить его координаты и определить объём хранения. Затем найти значение первого компонента во второй области.
\end{problem}

\begin{proof}[Решение]
Диадное произведение даёт матрицу
\[
F=sc^{\trans}
=
\begin{pmatrix}
2\\1
\end{pmatrix}
\begin{pmatrix}
0.5&0.3&0.2
\end{pmatrix}
=
\begin{pmatrix}
1&0.6&0.4\\
0.5&0.3&0.2
\end{pmatrix}.
\]
Строка соответствует пространственной области, столбец --- компоненту материала. Например,
\[
F_{21}=0.5
\]
есть значение первого компонента во второй области.

Полный массив содержит \(2\cdot3=6\) чисел. Разложимая запись хранит только \(2+3=5\) чисел: два пространственных коэффициента и три канальных. На больших массивах подобная низкоранговая структура может существенно уменьшать память и стоимость вычислений.
\end{proof}

\begin{problem}
\textit{Практическая задача: площадь элемента расчётной сетки.}
Треугольный элемент имеет вершины
\[
p_0=(1,1),
\qquad
p_1=(4,2),
\qquad
p_2=(2,5).
\]
С помощью внешнего произведения найти его площадь и определить ориентацию обхода \(p_0,p_1,p_2\).
\end{problem}

\begin{proof}[Решение]
Из вершины \(p_0\) выходят рёбра
\[
a=p_1-p_0=(3,1),
\qquad
b=p_2-p_0=(1,4).
\]
Их внешнее произведение равно
\[
a\wedge b
=
\det
\begin{pmatrix}
3&1\\
1&4
\end{pmatrix}
e_1\wedge e_2
=
11\,e_1\wedge e_2.
\]
Коэффициент \(11\) является ориентированной площадью параллелограмма. Поэтому площадь треугольника равна
\[
S=\frac12|11|=\frac{11}{2}.
\]
Знак положителен, следовательно, обход \(p_0,p_1,p_2\) имеет положительную, то есть против часовой стрелки, ориентацию. Для расчётной сетки ненулевое значение подтверждает, что элемент не вырожден.
\end{proof}

\section{Вычислительная линейная алгебра}
\label{sec:computational-linear-algebra}

Предыдущие разделы описывали точные свойства линейных объектов. В вычислениях дополнительно важно, насколько результат чувствителен к ошибкам данных и округления, каким способом хранить крупные матрицы и как получать устойчивые приближения при неполной или противоречивой информации.

Этот раздел расширяет материал Винберга темами, необходимыми для численных методов, анализа данных и генеративного проектирования. Основные конструкции --- сингулярное разложение, псевдообратная матрица, обусловленность, метод главных компонент и матричное дифференцирование.

\subsection{Матричные нормы и чувствительность вычислений}

Норма измеряет величину вектора или оператора. Для \(x\in\R^n\) наиболее часто используются
\[
\|x\|_1=\sum_{i=1}^{n}|x_i|,
\qquad
\|x\|_2=\left(\sum_{i=1}^{n}x_i^2\right)^{1/2},
\qquad
\|x\|_\infty=\max_i|x_i|.
\]

\begin{definition}
Операторной нормой матрицы \(A\), согласованной с выбранной векторной нормой, называется
\[
\|A\|
=
\max_{x\neq0}\frac{\|Ax\|}{\|x\|}
=
\max_{\|x\|=1}\|Ax\|.
\]
\end{definition}

Она показывает максимальный коэффициент усиления вектора линейным преобразованием. Для евклидовой нормы
\[
\|A\|_2=\sigma_{\max}(A),
\]
где \(\sigma_{\max}(A)\) --- наибольшее сингулярное число.

Для матрицы также используется норма Фробениуса
\[
\|A\|_F
=
\left(\sum_{i,j}a_{ij}^2\right)^{1/2}
=
\sqrt{\operatorname{tr}(A^{\trans}A)}.
\]
Она является евклидовой нормой массива коэффициентов и удобна при оптимизации по матрицам.

\begin{definition}
Пусть \(A\) обратима. Числом обусловленности матрицы в согласованной норме называется
\[
\kappa(A)=\|A\|\,\|A^{-1}\|.
\]
Для евклидовой нормы
\[
\kappa_2(A)
=
\frac{\sigma_{\max}(A)}{\sigma_{\min}(A)}.
\]
\end{definition}

Число обусловленности не описывает ошибку алгоритма. Оно описывает чувствительность самой задачи. Если \(\kappa(A)\) велико, то малое возмущение данных может вызвать большое относительное изменение решения независимо от выбранного корректного алгоритма.

Для системы
\[
Ax=b
\]
при малом возмущении правой части \(b+\delta b\) имеем
\[
A(x+\delta x)=b+\delta b,
\qquad
\delta x=A^{-1}\delta b.
\]
Отсюда следует оценка
\[
\frac{\|\delta x\|}{\|x\|}
\leq
\kappa(A)
\frac{\|\delta b\|}{\|b\|}.
\]
Она является верхней границей: фактическое усиление зависит от направления ошибки.

\begin{example}
Пусть
\[
A_\varepsilon=
\begin{pmatrix}
1&1\\
1&1+\varepsilon
\end{pmatrix}.
\]
При малом \(\varepsilon\) столбцы почти совпадают. Матрица формально обратима при \(\varepsilon\neq0\), но восстановление двух независимых параметров становится чувствительным к шуму. Плохая обусловленность здесь имеет содержательный смысл: два воздействия почти неразличимы по наблюдениям.
\end{example}

\subsection{Сингулярное разложение}

Собственные значения применимы непосредственно к квадратным операторам. Сингулярное разложение существует для любой прямоугольной матрицы и связывает её с двумя ортонормированными системами направлений.

\begin{theorem}[сингулярное разложение]
Для любой матрицы \(A\in\R^{m\times n}\) существуют ортогональные матрицы
\[
U\in\R^{m\times m},
\qquad
V\in\R^{n\times n}
\]
и прямоугольная диагональная матрица
\[
\Sigma\in\R^{m\times n}
\]
такие, что
\[
A=U\Sigma V^{\trans}.
\]
Диагональные элементы
\[
\sigma_1\geq\sigma_2\geq\cdots\geq\sigma_r>0
\]
называются сингулярными числами, а \(r=\operatorname{rank}A\).
\end{theorem}

Столбцы \(v_i\) матрицы \(V\) являются собственными векторами \(A^{\trans}A\):
\[
A^{\trans}Av_i=\sigma_i^2v_i.
\]
Столбцы \(u_i\) матрицы \(U\) являются собственными векторами \(AA^{\trans}\), а для ненулевых сингулярных чисел
\[
Av_i=\sigma_i u_i,
\qquad
A^{\trans}u_i=\sigma_i v_i.
\]

Геометрически преобразование выполняется в три этапа:
\[
x
\overset{V^{\trans}}{\longmapsto}
V^{\trans}x
\overset{\Sigma}{\longmapsto}
\Sigma V^{\trans}x
\overset{U}{\longmapsto}
U\Sigma V^{\trans}x.
\]
Сначала координаты поворачиваются в систему правых сингулярных направлений, затем независимо масштабируются коэффициентами \(\sigma_i\), после чего поворачиваются в пространстве результата.

Эквивалентная сумма имеет вид
\[
A
=
\sum_{i=1}^{r}\sigma_i u_i v_i^{\trans}.
\]
Каждое слагаемое имеет ранг один. Сингулярное разложение тем самым разбивает преобразование на независимые каналы передачи информации.

\subsection{Фундаментальные подпространства через SVD}

Пусть
\[
A=U\Sigma V^{\trans}
\]
и ненулевы ровно \(\sigma_1,\dots,\sigma_r\). Тогда
\[
\operatorname{im}A
=
\operatorname{span}(u_1,\dots,u_r),
\]
\[
\ker A
=
\operatorname{span}(v_{r+1},\dots,v_n),
\]
\[
\operatorname{im}A^{\trans}
=
\operatorname{span}(v_1,\dots,v_r),
\]
\[
\ker A^{\trans}
=
\operatorname{span}(u_{r+1},\dots,u_m).
\]

Таким образом, SVD одновременно показывает ранг, наблюдаемые направления, недостижимые выходы и невидимые изменения входа.

\subsection{Низкоранговое приближение}

Пусть
\[
A=\sum_{i=1}^{r}\sigma_i u_i v_i^{\trans}.
\]
Обрезанное разложение ранга \(k<r\) определяется формулой
\[
A_k
=
\sum_{i=1}^{k}\sigma_i u_i v_i^{\trans}.
\]

\begin{theorem}[Эккарта---Юнга]
Среди всех матриц \(B\) ранга не выше \(k\) матрица \(A_k\) минимизирует ошибки
\[
\|A-B\|_2
\quad\text{и}\quad
\|A-B\|_F.
\]
При этом
\[
\|A-A_k\|_2=\sigma_{k+1},
\]
\[
\|A-A_k\|_F^2
=
\sum_{i=k+1}^{r}\sigma_i^2.
\]
\end{theorem}

Теорема отвечает на практический вопрос: как сохранить только \(k\) наиболее значимых линейных режимов, потеряв как можно меньше информации в выбранной норме.

Низкоранговое приближение применяется при сжатии данных, удалении шума, ускорении линейных слоёв и построении компактных пространств параметров формы. Значение \(k\) является модельным выбором: слишком малое значение уничтожает существенные направления, слишком большое сохраняет шум и избыточность.

\subsection{Псевдообратная матрица}

Для прямоугольной или вырожденной матрицы обычной обратной матрицы нет. Тем не менее можно выбрать решение, которое наилучшим образом согласуется с данными и имеет минимальную норму.

Пусть
\[
A=U\Sigma V^{\trans}.
\]
Матрица \(\Sigma^+\) получается заменой каждого ненулевого \(\sigma_i\) на \(1/\sigma_i\) и транспонированием прямоугольной диагональной формы.

\begin{definition}
Псевдообратной матрицей Мура---Пенроуза называется
\[
A^+=V\Sigma^+U^{\trans}.
\]
\end{definition}

Она единственна и удовлетворяет условиям
\[
AA^+A=A,
\qquad
A^+AA^+=A^+,
\]
\[
(AA^+)^{\trans}=AA^+,
\qquad
(A^+A)^{\trans}=A^+A.
\]

Матрица
\[
AA^+
\]
является ортогональным проектором на \(\operatorname{im}A\), а
\[
A^+A
\]
--- ортогональным проектором на \(\operatorname{im}A^{\trans}\).

Для любой системы \(Ax=b\) вектор
\[
x_*=A^+b
\]
минимизирует невязку
\[
\|Ax-b\|_2.
\]
Среди всех решений с минимальной невязкой он имеет наименьшую норму \(\|x\|_2\).

Если столбцы \(A\) линейно независимы, то
\[
A^+=(A^{\trans}A)^{-1}A^{\trans}.
\]
Если строки \(A\) линейно независимы, то
\[
A^+=A^{\trans}(AA^{\trans})^{-1}.
\]

\subsection{Регуляризация плохо обусловленных задач}

В формуле псевдообратной матрицы направления с малыми \(\sigma_i\) умножаются на большие коэффициенты \(1/\sigma_i\). Поэтому шум в таких направлениях резко усиливается.

Один способ стабилизации --- отбросить сингулярные числа ниже выбранного порога. Другой способ --- регуляризация Тихонова:
\[
x_\lambda
=
\arg\min_x
\left(
\|Ax-b\|_2^2+\lambda\|x\|_2^2
\right),
\qquad
\lambda>0.
\]
Условие стационарности даёт
\[
(A^{\trans}A+\lambda I)x_\lambda
=
A^{\trans}b.
\]
Следовательно,
\[
x_\lambda
=
(A^{\trans}A+\lambda I)^{-1}A^{\trans}b.
\]

В сингулярных координатах коэффициент \(1/\sigma_i\) заменяется на
\[
\frac{\sigma_i}{\sigma_i^2+\lambda}.
\]
При малом \(\sigma_i\) этот множитель остаётся ограниченным. Цена устойчивости состоит в смещении решения к нулю. Параметр \(\lambda\) определяет компромисс между соответствием данным и устойчивостью.

\subsection{Метод главных компонент}

Пусть имеются наблюдения
\[
x_1,\dots,x_N\in\R^d.
\]
Сначала вычисляется среднее
\[
\bar x
=
\frac1N\sum_{i=1}^{N}x_i,
\]
после чего формируется центрированная матрица данных
\[
X=
\begin{pmatrix}
(x_1-\bar x)^{\trans}\\
\vdots\\
(x_N-\bar x)^{\trans}
\end{pmatrix}
\in\R^{N\times d}.
\]

Ковариационная матрица имеет вид
\[
C=\frac1N X^{\trans}X.
\]
Она симметрична и неотрицательно определена.

\begin{definition}
Первой главной компонентой называется единичный вектор \(v_1\), максимизирующий дисперсию проекций:
\[
v_1
=
\arg\max_{\|v\|_2=1}
\frac1N\|Xv\|_2^2.
\]
\end{definition}

Поскольку
\[
\frac1N\|Xv\|_2^2
=
v^{\trans}Cv,
\]
первая компонента является собственным вектором \(C\), соответствующим наибольшему собственному значению. Остальные компоненты выбираются последовательно при условии ортогональности.

Если
\[
X=U\Sigma V^{\trans},
\]
то столбцы \(V\) являются направлениями главных компонент, а
\[
\lambda_i(C)=\frac{\sigma_i^2}{N}.
\]

Проекция наблюдения \(x\) на первые \(k\) компонент равна
\[
z=V_k^{\trans}(x-\bar x),
\]
а приближённое восстановление ---
\[
\widehat x
=
\bar x+V_kz.
\]

Доля объяснённой дисперсии первых \(k\) компонент вычисляется как
\[
\frac{\sum_{i=1}^{k}\sigma_i^2}
{\sum_{i=1}^{r}\sigma_i^2}.
\]
PCA является линейным методом: он обнаруживает только линейные направления изменчивости и зависит от масштаба исходных признаков.

\subsection{Матричное дифференцирование}

Для оптимизации необходимо вычислять производные функций, аргументами которых являются векторы и матрицы. Удобнее всего начинать с дифференциала.

Пусть
\[
f:\R^n\to\R.
\]
Дифференциал записывается как
\[
df
=
\nabla f(x)^{\trans}dx.
\]
Коэффициент при \(dx\) определяет градиент.

Для функции матрицы
\[
f:\R^{m\times n}\to\R
\]
используется скалярное произведение Фробениуса:
\[
df
=
\langle\nabla_X f,dX\rangle_F
=
\operatorname{tr}\left(
(\nabla_X f)^{\trans}dX
\right).
\]

Основные правила следуют из обычной линейности:
\[
d(Ax)=A\,dx
\]
для постоянной матрицы \(A\),
\[
d(XY)=dX\,Y+X\,dY,
\]
\[
d(X^{\trans})=(dX)^{\trans}.
\]

Для квадратичной функции
\[
f(x)=\frac12x^{\trans}Ax
\]
имеем
\[
df
=
\frac12(dx)^{\trans}Ax
+
\frac12x^{\trans}A\,dx.
\]
Приводя оба слагаемых к форме коэффициента при \(dx\), получаем
\[
\nabla f(x)
=
\frac12(A+A^{\trans})x.
\]
Если \(A\) симметрична, то
\[
\nabla f(x)=Ax.
\]

Для функции наименьших квадратов
\[
L(x)=\frac12\|Ax-b\|_2^2
\]
обозначим \(r=Ax-b\). Тогда
\[
dL
=
r^{\trans}A\,dx,
\]
поэтому
\[
\nabla L(x)
=
A^{\trans}(Ax-b).
\]
Условие \(\nabla L(x)=0\) снова приводит к нормальным уравнениям
\[
A^{\trans}Ax=A^{\trans}b.
\]

Для матричной задачи
\[
F(X)
=
\frac12\|AX-B\|_F^2
\]
аналогично получается
\[
\nabla_XF
=
A^{\trans}(AX-B).
\]

\subsection{След и полезные производные}

След позволяет циклически переставлять множители:
\[
\operatorname{tr}(ABC)
=
\operatorname{tr}(BCA)
=
\operatorname{tr}(CAB),
\]
если размеры матриц согласованы. Это свойство используется для переноса \(dX\) в конец выражения.

Для постоянных матриц подходящих размеров:
\[
\nabla_X\operatorname{tr}(A^{\trans}X)=A,
\]
\[
\nabla_X\frac12\|X\|_F^2=X,
\]
\[
\nabla_X\frac12\|AXB-C\|_F^2
=
A^{\trans}(AXB-C)B^{\trans}.
\]

Если \(X\) обратима, то из
\[
XX^{-1}=I
\]
следует
\[
d(X^{-1})
=
-X^{-1}(dX)X^{-1}.
\]
Эта формула нужна при дифференцировании решений линейных систем и ковариационных моделей.

\paragraph{Связь с обратным распространением.}

Обратное распространение ошибки является систематическим применением правила цепочки к композиции отображений. Для линейного слоя
\[
y=Wx+b
\]
и скалярной функции потерь \(L(y)\) обозначим
\[
g=\nabla_yL.
\]
Тогда
\[
\nabla_xL=W^{\trans}g,
\]
\[
\nabla_WL=gx^{\trans},
\qquad
\nabla_bL=g.
\]

Транспонированная матрица появляется не как условность программной библиотеки, а как действие сопряжённого линейного отображения на функционал изменения потерь.

\paragraph{Связь с генеративным проектированием.}

SVD и PCA выделяют главные линейные режимы изменения формы, поля или набора параметров. Псевдообратная матрица восстанавливает управляющие воздействия по избыточным или неполным измерениям. Регуляризация подавляет неустойчивые направления. Матричное дифференцирование позволяет оптимизировать параметры геометрии через симулятор или дифференцируемую модель.

Эти методы не отменяют содержательной постановки задачи. Малое сингулярное число может означать шум, но может также соответствовать редкому физически важному режиму. Главная компонента описывает наибольшую изменчивость, а не обязательно наибольшую полезность. Поэтому автоматическое усечение направлений должно сопровождаться физической интерпретацией.

\subsection*{Задачи}

\begin{problem}
\textit{Задача по материалу Винберга.}
Пусть \(A\in\R^{m\times n}\) имеет линейно независимые столбцы. Показать, что
\[
P=A(A^{\trans}A)^{-1}A^{\trans}
\]
является ортогональным проектором на \(\operatorname{im}A\), а
\[
A^+=(A^{\trans}A)^{-1}A^{\trans}.
\]
\end{problem}

\begin{proof}[Решение]
Матрица \(A^{\trans}A\) обратима, потому что
\[
x^{\trans}A^{\trans}Ax
=
\|Ax\|_2^2
\]
равно нулю только при \(x=0\).

Проверим симметричность:
\[
P^{\trans}
=
A(A^{\trans}A)^{-1}A^{\trans}
=
P.
\]
Далее
\[
P^2
=
A(A^{\trans}A)^{-1}
A^{\trans}A
(A^{\trans}A)^{-1}A^{\trans}
=
P.
\]
Следовательно, \(P\) является ортогональным проектором.

Для любого \(y=Ax\in\operatorname{im}A\)
\[
Py
=
A(A^{\trans}A)^{-1}A^{\trans}Ax
=
Ax=y.
\]
Кроме того, результат \(Py\) всегда принадлежит \(\operatorname{im}A\). Значит, проектирование выполняется именно на образ \(A\).

Матрица
\[
A^+=(A^{\trans}A)^{-1}A^{\trans}
\]
задаёт коэффициенты проекции:
\[
AA^+b=Pb.
\]
Поэтому \(x_*=A^+b\) является решением задачи наименьших квадратов.
\end{proof}

\begin{problem}
\textit{Задача по материалу Винберга.}
Для матрицы
\[
A=
\begin{pmatrix}
3&0\\
0&1
\end{pmatrix}
\]
найти лучшее приближение ранга один в норме Фробениуса и вычислить ошибку.
\end{problem}

\begin{proof}[Решение]
Матрица уже имеет сингулярное разложение
\[
A
=
3e_1e_1^{\trans}
+
1e_2e_2^{\trans}.
\]
Сингулярные числа равны
\[
\sigma_1=3,
\qquad
\sigma_2=1.
\]
По теореме Эккарта---Юнга лучшее приближение ранга один получается отбрасыванием второго слагаемого:
\[
A_1
=
\begin{pmatrix}
3&0\\
0&0
\end{pmatrix}.
\]
Ошибка равна
\[
\|A-A_1\|_F
=
\left\|
\begin{pmatrix}
0&0\\
0&1
\end{pmatrix}
\right\|_F
=1.
\]
Сохраняется направление, которое преобразование усиливает в три раза, а более слабое направление отбрасывается.
\end{proof}

\begin{problem}
\textit{Практическая задача: калибровка двух почти одинаковых приводов.}
Два привода изменяют положения двух контрольных точек. Линеаризованная модель имеет вид
\[
A=
\begin{pmatrix}
1&1\\
1&1.01
\end{pmatrix},
\qquad
b=
\begin{pmatrix}
2\\
2.01
\end{pmatrix}.
\]
Найти точное решение. Затем объяснить, почему малые ошибки измерения могут сильно изменить найденные команды и как регуляризация исправляет ситуацию.
\end{problem}

\begin{proof}[Решение]
Из первого уравнения
\[
x_1+x_2=2.
\]
Вычитая его из второго,
\[
0.01x_2=0.01,
\]
поэтому
\[
x_2=1,
\qquad
x_1=1.
\]

Однако столбцы матрицы почти совпадают:
\[
\begin{pmatrix}1\\1\end{pmatrix},
\qquad
\begin{pmatrix}1\\1.01\end{pmatrix}.
\]
Система различает перераспределение команды между приводами только по разности \(0.01\). Ошибка второй координаты \(b\) порядка \(10^{-2}\) может изменить \(x_2\) на величину порядка единицы.

Регуляризованное решение определяется системой
\[
(A^{\trans}A+\lambda I)x_\lambda
=
A^{\trans}b.
\]
Добавка \(\lambda I\) запрещает чрезмерно большие и взаимно компенсирующие команды. При этом решение перестаёт точно воспроизводить измерения, но становится менее чувствительным к шуму. Физически регуляризация выражает предпочтение умеренных управляющих воздействий.
\end{proof}

\begin{problem}
\textit{Практическая задача: главные режимы формы.}
Три центрированных измерения двух геометрических параметров имеют вид
\[
x_1=(-1,-1),
\qquad
x_2=(0,0),
\qquad
x_3=(1,1).
\]
Найти первую главную компоненту, спроецировать на неё наблюдения и объяснить геометрический смысл результата.
\end{problem}

\begin{proof}[Решение]
Матрица данных равна
\[
X=
\begin{pmatrix}
-1&-1\\
0&0\\
1&1
\end{pmatrix}.
\]
Ковариационная матрица с множителем \(1/3\) имеет вид
\[
C
=
\frac13X^{\trans}X
=
\frac13
\begin{pmatrix}
2&2\\
2&2
\end{pmatrix}.
\]
Её собственные направления:
\[
v_1
=
\frac1{\sqrt2}
\begin{pmatrix}
1\\1
\end{pmatrix},
\qquad
v_2
=
\frac1{\sqrt2}
\begin{pmatrix}
1\\-1
\end{pmatrix}.
\]
Первому направлению соответствует ненулевое собственное значение, второму --- нулевое. Поэтому первая главная компонента равна \(v_1\).

Координаты наблюдений:
\[
z_i=v_1^{\trans}x_i,
\]
то есть
\[
z_1=-\sqrt2,
\qquad
z_2=0,
\qquad
z_3=\sqrt2.
\]
Все данные точно лежат на одной прямой. Два геометрических параметра изменяются совместно и могут быть заменены одной координатой формы без потери информации в данном наборе.
\end{proof}

% HIDDENPOWER FINAL STYLE AUDIT APPLIED
